% ======================================================================
%  PAPER I of the two-paper split -- JMP VARIANT
%  Controlled-rail brachistochrones in non-stationary spacetimes:
%  conformal symmetry, Kodama energy, and Vaidya dynamics
%  Target: Journal of Mathematical Physics (AIP Publishing)
%  (spherically symmetric sector: FLRW + Vaidya; shared control-theory
%   foundations live here and are cited by Paper II)
%
%  Difference from paper1.tex: the six shared results I.1-I.6 are cast as
%  genuine amsthm environments with explicit hypotheses (H1)-(H3) and
%  delimited proofs, plus a main-results roadmap and the AIP-required
%  author declarations. Physics content is unchanged.
% ======================================================================
\documentclass[12pt]{iopart}

% iopart pre-defines \equation*, which clashes with amsmath; free it first
\makeatletter
\expandafter\let\csname equation*\endcsname\relax
\expandafter\let\csname endequation*\endcsname\relax
\makeatother
\usepackage{amsmath}
\usepackage{amssymb}
\usepackage{amsthm}
\usepackage{iopams}
\usepackage{graphicx}
\usepackage{array}
\usepackage{bm}
\usepackage{xcolor}
\usepackage{placeins}
\usepackage{hyperref}
\hypersetup{
  colorlinks=true,
  linkcolor=blue,
  citecolor=blue,
  urlcolor=blue,
  pdftitle={Controlled-rail brachistochrones in non-stationary spacetimes: conformal symmetry, Kodama energy, and Vaidya dynamics},
  pdfauthor={Iman Rosignoli},
  pdfsubject={Journal of Mathematical Physics submission},
  pdfkeywords={relativistic brachistochrone; controlled rail; Kodama energy; Vaidya spacetime; FLRW; optimal control; Pontryagin maximum principle; Hamilton-Jacobi verification; Randers-Zermelo geometry; abelian integrals}
}

% figures shared with the master/Paper II
\graphicspath{{../paper/}}

% prose placeholder helper
\newcommand{\todo}[1]{{\color{red}\small[[#1]]}}

% shortcuts
\newcommand{\Ehat}{\hat E}
\newcommand{\rp}{r_{+}}
\newcommand{\re}{r_{e}}
\newcommand{\Jc}{J_{\rm deg}}
\newcommand{\Jpen}{J_{\rm pen}}
\newcommand{\Af}{A_{\mathrm{freeze}}}
\newcommand{\dd}{\mathrm{d}}

% ----------------------------------------------------------------------
%  Theorem environments.  The six shared results carry the fixed labels
%  I.1-I.6 used by the companion Paper II, so they are NOT auto-numbered
%  in order of appearance (I.4 and I.6 live in the appendices while I.5
%  is stated in section 4).  Each environment is preceded by an explicit
%  \setcounter so that the printed number is stable under any reordering.
% ----------------------------------------------------------------------
\theoremstyle{plain}
\newtheorem{theorem}{Theorem}
\newtheorem{proposition}[theorem]{Proposition}
\newtheorem{lemma}[theorem]{Lemma}
\newtheorem{corollary}[theorem]{Corollary}
\renewcommand{\thetheorem}{I.\arabic{theorem}}
\theoremstyle{definition}
\newtheorem{hypothesis}{Hypothesis}
\renewcommand{\thehypothesis}{H\arabic{hypothesis}}
% auxiliary results introduced by the present variant carry letters, so that
% the numbers I.1-I.6 cited by the companion paper are never disturbed
\theoremstyle{plain}
\newtheorem{auxlemma}{Lemma}
\newtheorem{auxtheorem}[auxlemma]{Theorem}
\renewcommand{\theauxlemma}{I.\Alph{auxlemma}}
% Remarks are deliberately unnumbered: they must not consume the shared
% theorem counter, whose values are pinned to the stable labels I.1-I.6.
\theoremstyle{remark}
\newtheorem*{remark}{Remark}
% \Ithm{k} fixes the number of the next theorem-like environment to I.k
\newcommand{\Ithm}[1]{\setcounter{theorem}{\numexpr#1-1\relax}}

\begin{document}

\title{Controlled-rail brachistochrones in non-stationary spacetimes:
conformal symmetry, Kodama energy, and Vaidya dynamics}

\author{Iman Rosignoli}
\address{Department of Physics, University of Pavia, Via A.\ Bassi 6, 27100 Pavia, Italy}
\ead{iman.rosignoli01@universitadipavia.it}

\begin{abstract}
What is the fastest constrained worldline in a spacetime that is itself evolving? In a
stationary spacetime the relativistic brachistochrone reduces to Fermat/Randers optics
with a conserved rail energy supplied by a timelike Killing vector; in a dynamical
spacetime no such conservation law exists and the variational problem becomes
non-autonomous. We address this by formulating the brachistochrone as a
\emph{controlled-rail} optimal-control problem, in which the invariant $-u\cdot W=\Ehat$ is
actively maintained along the worldline by a selector $W$ that follows a Killing $\to$
conformal-Killing $\to$ Kodama hierarchy and reduces to the \emph{Kodama} vector when no
timelike Killing field survives. We show that this construction is a legitimate
Pontryagin problem---establishing existence and normality on a regular timelike-selector
domain and a non-autonomous Hamilton--Jacobi \emph{verification} criterion (global
minimisation being conditional, not automatic)---and then derive the extremal equations
and their closed representations in two spherically symmetric non-stationary cases:
Friedmann--Lema\^itre--Robertson--Walker (FLRW) at arbitrary spatial curvature and the
ingoing Vaidya dynamical black hole. The domain is delimited at the outset:
$g(K,K)=-(1-2m(v)/r)$ for the Kodama selector, so the admissible-velocity set is
compact only for $r>2m(v)$ and the compact-control problem \emph{ends} at $r=2m(v)$;
what the Vaidya calculation establishes is an exterior approach and contact threshold
rather than an interior optimal trajectory, and $m'(v)<0$ in the ingoing metric is a
formal continuation, physical evaporation requiring the outgoing problem. For Vaidya we
obtain the approach phenomenology and a complete
first-order adiabatic correction, verified against the true non-autonomous flow, that
separates the homogeneous-expansion effects of the FLRW base from the radial
spatial-gradient and mass-flow effects of the dynamical horizon. The rotating,
axisymmetric conformal-Kerr (Thakurta--Kerr)
application is developed in a companion paper.
\end{abstract}

\maketitle
% Table of contents removed for the research-paper submission (the running
% header no longer displays CONTENTS; the journal reformats the accepted article).
% Concise running head to avoid truncation of the long title.
\markboth{Controlled-rail brachistochrones in non-stationary spacetimes}%
         {Controlled-rail brachistochrones in non-stationary spacetimes}

% ======================================================================
\section{Introduction}
\label{sec:intro}
The brachistochrone---the path of fastest descent---is a classical variational problem
whose relativistic analogue asks for the fastest constrained ``rail'' worldline from a
launch event to a fixed spatial target, at a \emph{free} value of the arrival clock. In a
\emph{stationary} spacetime the answer is well understood: a timelike Killing vector
supplies a conserved rail energy, and the problem reduces to Fermat's principle and to
optical (Randers--Zermelo--Finsler) geometry
\cite{Perlick1991}, \cite[\S\S6.5--6.6,~7.3]{PerlickRayOptics},
\cite{Randers1941,GibbonsWerner2008,GibbonsRanders2009,BaoRoblesShen}---the first
for the timelike brachistochrone itself, the monograph for the ray-optical and
Fermat machinery it uses.
Existence, multiplicity and Morse theory for these stationary arrival- and travel-time
brachistochrones, together with their sub-Riemannian and Fermat-principle formulations,
are by now well developed
\cite{GiannoniPiccioneVerderesi1997,GiannoniPiccioneTausk2002,GiannoniPiccione2002,CaponioJavaloyes2026}.
The physically central regimes of gravitation, however---cosmological expansion,
gravitational collapse, accretion---are \emph{not} stationary.

\paragraph{Where this problem parts company with the established one.} The
relation to Giannoni and Piccione's arrival-time brachistochrone
\cite{GiannoniPiccione2002} is worth stating exactly, because it locates the
present contribution more sharply than a novelty claim can. Their trial paths are
timelike, as ours are, and they constrain the four-acceleration to be perpendicular
to the plane spanned by $\dot\sigma$ and the observer field $Y$---a
\emph{frictionless slide}, in which the constraint force does no work against $Y$
\cite[Eqs.~(1.5)--(1.6)]{GiannoniPiccione2002}. Differentiating the charge along
the curve gives, for any $W$ and any timelike $u$,
\begin{equation}
    \frac{\dd}{\dd\tau}\big(\!-g(u,W)\big)
      = -\,g(a,W)-\tfrac12(\mathcal L_Wg)_{\mu\nu}u^\mu u^\nu ,
    \label{eq:gp-fork}
\end{equation}
so under their condition $g(a,W)=0$ the charge drifts at exactly the Killing
defect---the same rate as for a free geodesic. Their constraint force cannot hold
it. The two problems therefore coincide precisely when $W$ is Killing, where the
defect vanishes and their energy relation $g(\dot\sigma,Y)=-k\mathcal T_\sigma$
\cite[Eq.~(1.7)]{GiannoniPiccione2002} \emph{is} the rail invariant; and they
diverge as soon as it is not.

That divergence has a consequence the same authors record. They observe that in the
non-stationary case it is ``not even clear whether the non-stationary
brachistochrones are solutions to a second-order differential equation'', that the
Lagrange multipliers cannot be eliminated ``unless in the stationary case'', and
that the brachistochrones there appear to be governed by an integro-differential
rather than a differential equation \cite[\S1]{GiannoniPiccione2002}. The
controlled rail is the other branch of~\eqref{eq:gp-fork}: instead of keeping
$g(a,W)=0$ and losing the charge, it keeps the charge and pays
$g(a,W)=-\tfrac12(\mathcal L_Wg)_{\mu\nu}u^\mu u^\nu$, which is the thrust. That
substitution is what leaves the non-stationary problem a second-order control
system with a compact control set, and it is the reason a maximum principle applies
here at all. Equation~\eqref{eq:gp-fork} is verified in
\texttt{verification/verify\_gp2002\_fork.wls}.

A word is owed to Kovner's Fermat principle \cite{Kovner1990}, which is not confined
to the stationary case and must therefore be distinguished carefully rather than
bundled with it. He states his principle for ``arbitrary, stationary, and
nonstationary metrics, for massive and massless particles'', in the form: an
observer is intercepted by \emph{geodesics} which are extremals of the observer's
worldline parameter with respect to path variations \emph{restricted to the mass
shell} \cite[Eq.~(2.14)]{Kovner1990}. Three things separate that from the present
problem, and none of them is a matter of degree. His extremals are geodesics---his
principle is a characterisation of free fall---whereas the controlled rail is
\emph{not} a geodesic: it carries the proper acceleration
$a\cdot W=-u^au^b\nabla_{(a}W_{b)}$ that maintains the invariant, and vanishing
thrust is exactly the case his principle covers. His only restriction on the
competing paths is the mass shell; ours is the mass shell \emph{together with}
$-u\cdot W=\Ehat$, and it is that second constraint which turns the velocity locus
into an oval and the problem into one of optimal control. He asks for
\emph{extremality}, noting that the arrival time at such an extremal may be a
minimum, a maximum or a saddle, whereas the question pursued here is genuine
minimality. And his non-stationary content is perturbative: the one genuinely
time-dependent application is a \emph{static} lens perturbed by a weak
gravitational wave, worked to first order to obtain the transverse velocity of
caustic motion and thence to test an astrophysical proposal about $\gamma$-ray
bursts \cite[\S\S IV--V]{Kovner1990}. The question there is where images form and
when they arrive, not what the time-optimal constrained worldline is. So Kovner's principle applies to non-stationary spacetimes and to
massive particles, and it does not address the controlled-rail problem; what is new
here is not a Fermat principle for evolving metrics but the maintenance of a
prescribed invariant along a non-geodesic worldline, and the optimal-control theory
that follows from it.

These treatments all exploit stationarity---a prescribed observer-field energy, a
conserved Killing charge, or a time-independent optical metric. When the background
geometry itself evolves that structure disappears: there is no timelike Killing field, no
conserved energy to label the extremals, and the reduced problem is genuinely
non-autonomous, its ordinary Hamiltonian no longer an interior first integral. The
difficulty we address is therefore specific: how to \emph{prescribe and maintain} a rail
charge in an evolving spacetime, what field should supply it, and how time
dependence---as opposed to spatial gradients---reshapes the optimal path.

Our starting point is that the rail energy need not be a symmetry charge: it can be an
invariant that is \emph{actively maintained} along the worldline. We formulate the
brachistochrone as a \emph{controlled-rail} optimal-control problem in which
$-u\cdot W=\Ehat$ is held fixed by a selector $W$ following a Killing $\to$
conformal-Killing $\to$ Kodama hierarchy; in a dynamical spherically symmetric spacetime
$W$ is the \emph{Kodama} vector, whose charge is conserved along the controlled worldline
even where free fall is not. Pontryagin's maximum principle makes this a legitimate
time-optimal control problem, supplying existence, normality and a non-autonomous
Hamilton--Jacobi verification criterion in place of the stationary conservation law.

We develop this general framework in section~\ref{sec:principle} and then apply it to two
spherically symmetric non-stationary spacetimes: the spatially flat FLRW base
(section~\ref{sec:flrw}) and the ingoing Vaidya dynamical black hole
(section~\ref{sec:vaidya}). For Vaidya we obtain the plunge phenomenology and a
first-order adiabatic correction in closed special-function form, verified against the
true non-autonomous flow; the construction cleanly separates homogeneous expansion from
radial spatial-gradient and mass-flow effects. The rotating, axisymmetric conformal-Kerr
(Thakurta--Kerr) case---frame dragging, an ergosphere and genus-two closed forms---is
developed in the companion Paper~II~\cite{PaperTwo}, which cites the foundations proved
here.

\paragraph{Main results and organization.} Six results carry the mathematical content of
the paper and are labelled I.1--I.6 throughout, a numbering that the companion paper cites
verbatim. All six are stated under the standing hypotheses (H1)--(H3) of
section~\ref{sec:indicatrix}, which delimit the \emph{regular domain} by three checkable
scalar inequalities: the selector timelike, the rail energy strictly above the local
freezing value $|W|$, and that margin uniform on a compact domain. The geometric
properties actually used---that the velocity indicatrix is a compact strictly convex
ellipse and that the arrival-time functional is two-sidedly bounded---are not assumed but
\emph{derived} from those inequalities in Lemma~\ref{lem:indicatrix}, which also identifies
the freezing surface as the locus $\Ehat=|W|$ where the indicatrix collapses to a point.
\begin{itemize}\itemsep2pt
\item \textbf{Theorem~\ref{thm:rail}} (section~\ref{sec:indicatrix}) reduces the
controlled-rail brachistochrone to a time-optimal control problem with compact control set
$S^1$, whose maximized Hamiltonian is single-valued in closed form and requires no
Filippov relaxation.
\item \textbf{Theorem~\ref{thm:existence}} (section~\ref{sec:control-foundations})
establishes existence of a minimizer by the direct method and the \emph{normality} of every
regular extremal, so that the reduction to the rail invariant is never vacuous.
\item \textbf{Proposition~\ref{prop:hjb}} (section~\ref{sec:control-foundations}) supplies
the non-autonomous Hamilton--Jacobi verification criterion that upgrades a saturated
extremal to a global minimizer, together with an explicit statement of what is
\emph{not} certified.
\item \textbf{Lemma~\ref{lem:reduction}} (\ref{app:reproduce}) shows that the whole
adiabatic hierarchy collapses onto a single degree-ten polynomial identity, hence a
nonsingular $11\times11$ linear system whose solution is rational in the metric
parameters: every coefficient in the paper is solved for, none is fitted.
\item \textbf{Theorem~\ref{thm:firstorder}} (section~\ref{sec:vaidya-adiabatic}, proved in
\ref{app:firstorder}) gives the complete first-order response of the extremal to the slow
drift---on-shell \emph{and} off-shell---in a metric-independent form whose source collapses
to the universal boundary term $S_{\mathcal D}=[\,r\,p_r\,]-\lambda$.
\item \textbf{Lemma~\ref{lem:tracking}} (\ref{app:tracking}) computes the derivative of the
genus-degeneration family along the drifting locus and shows that it annihilates the
leading (triple) pole of the source.
\end{itemize}
The remaining sections are applications of this apparatus: section~\ref{sec:flrw} treats
the degenerate FLRW base, section~\ref{sec:vaidya} the ingoing Vaidya black hole, where
the closed forms are abelian integrals on a genus-two curve completed by a horizon
dilogarithm. Appendix~\ref{app:validation} records the reproducibility and consistency
data, \ref{app:degeneration} the explicit forms on the genus-degeneration loci, and
\ref{app:firstorder} the proof of Theorem~\ref{thm:firstorder}.

\section{The controlled-rail principle}
\label{sec:principle}
% ----------------------------------------------------------------------

\subsection{Rail constraint and the indicatrix}
\label{sec:indicatrix}
The rail is a constrained worldline with the invariant
\begin{equation}
    -\,u\cdot W \;=\; \Ehat ,
    \label{eq:rail}
\end{equation}
actively maintained by an ideal rail force $f_\mu=Du_\mu/\dd\tau$ with
$f\cdot u=0$ (automatic for a rest-mass--preserving force), which creates no rest
mass but is free to exert power against the selector (see below).

\begin{table}[htbp]
    \centering\small
    \caption{Notation used throughout. Rail charges $(\Ehat,J)$ are the conserved control
    costates; rotating conformal normalizations are outside the scope of this paper.}
    \label{tab:notation}
    \begin{tabular}{@{}ll@{}}
        \hline\hline
        symbol & meaning\\
        \hline
        $\Ehat=-u\cdot W$ & rail energy invariant (conserved control charge)\\
        $J=p_\varphi$ & axial control costate (conserved; \emph{not} $g(u,\partial_\varphi)$)\\
        $L_{\rm mech}=g(u,\partial_\varphi)$ & mechanical angular momentum (\emph{not} conserved on the rail)\\
        $m(v)$ & Vaidya Misner--Sharp mass\\
        $M$ & fixed reference mass used for dimensionless Vaidya variables\\
        $t,\tau,v$ & coordinate, proper, advanced-null clocks\\
        $\chi$ & monotone background time coordinate included in the augmented state\\
        $\mathcal T$ & selected cost clock; $\ell_{\mathcal T}=\dd\mathcal T/\dd\chi$\\
        $\sigma$ & arbitrary monotone parameter used to describe a spatial arc\\
        $s=v/M$ & dimensionless slow coordinate for the Vaidya expansion\\
        $\lambda$ & elapsed clock from the free-arrival endpoint\\
        $\dot m=\dd m/\dd v$ & Vaidya mass rate\\
        $\varepsilon=M\dot m/m$ & dimensionless Vaidya adiabatic parameter\\
        \hline\hline
    \end{tabular}
\end{table}

\paragraph{The ellipse.} Fix an event and a time function adapted to $W$. The
linear constraint~\eqref{eq:rail} fixes the $W$-component of $u$; the mass shell
$u\cdot u=-1$ is a quadric; their intersection, projected onto the physical
velocity plane $(\dd r,\,r\,\dd\varphi)$ per unit selector-time, is an ellipse
\begin{equation}
    \dot x(\theta)=c+R(\theta),\qquad \theta\in[0,2\pi),
    \label{eq:ellipse}
\end{equation}
whose center $c$ is the Zermelo wind (zero for FLRW, ingoing radial for Vaidya,
azimuthal---frame dragging---for Thakurta--Kerr) and whose axes $R$ shrink with
the local speed as the freezing surface is approached (figure~\ref{fig:indicatrici}).
Every cost clock---arrival/coordinate time $t$, proper time $\tau$, conformal time
$\eta$---shares this \emph{same} indicatrix; the branches differ only in the cost
one-form, i.e.\ in which clock measures the elapsed travel time. Two of these are
\emph{not} independent spatial families: since $\dd t=A(\eta)\,\dd\eta$ with $A>0$ is
a monotone reparametrization, the coordinate-time and conformal-time functionals are
extremized by the \emph{same} spatial path ($\arg\min t_f=\arg\min\eta_f$), so $\eta$
is the canonical conformal-time \emph{representation} of the arrival branch rather
than a third independent brachistochrone. The genuinely distinct branch is the proper
time $\tau$, whose cost carries the extra $\sqrt{1-v^2}$ weight; for Vaidya the native
arrival clock is the advanced null time $v$, again distinct from $\tau$.

\paragraph{Domain of the indicatrix.} The intersection~\eqref{eq:ellipse} is a
compact, strictly convex oval---the hypothesis under which the Pontryagin reduction
below applies---precisely where the selector $W$ is \emph{timelike} and the quadric
intersection is nondegenerate. Three loci require separate treatment. (i) At the
freezing surface the entire indicatrix collapses to a point, $R\to0$, and strict
convexity is lost; this degeneration is distinct from a radial turning point and
from a spectral separatrix, both treated separately below. (ii) At the horizon the
radial metric coefficient degenerates. (iii) For the rotating conformal
Thakurta--Kerr case, inside the \emph{ergosphere} $\partial_\eta$---hence
$W$---becomes spacelike and the mass-shell intersection is no longer automatically a
compact oval in the PMP sense; that regime, and the alternative timelike-congruence
re-anchoring it would require, is developed in the companion Paper~II. The local
theorems of this section are therefore stated in the timelike-$W$, nondegenerate
domain.

\paragraph{The optimal-control problem.} Precisely, let $x=(r,\varphi)$ denote the
spatial position and let $\chi$ be a monotone background time coordinate ($\chi=\eta$
for FLRW and $\chi=v$ for ingoing Vaidya). The full non-autonomous state is therefore
$y=(\chi,x)$, even when only its spatial projection is displayed. An arbitrary monotone
curve parameter is denoted by $\sigma$. The control is the unit spatial direction
$\theta\in S^1$ on the rail indicatrix~\eqref{eq:ellipse}; using $\chi$ as parameter gives
$\dd x/\dd\chi=f(x,\theta,\chi)$. A chosen clock $\mathcal T$ is the \emph{cost},
\[
    J_{\mathcal T}=\int \ell_{\mathcal T}(x,\theta,\chi)\,\dd\chi,
    \qquad \ell_{\mathcal T}=\frac{\dd\mathcal T}{\dd\chi}.
\]
Thus $\ell_{\mathcal T}=1$ when the cost clock itself is the evolution coordinate,
whereas for the Vaidya proper-time branch $\chi=v$ and
$\ell_\tau=\dd\tau/\dd v$. The boundary data are a \emph{fixed initial event}
$(\chi_0,x_0)$ and a \emph{fixed spatial target} $x_f$ (a point, or a target
set/worldline), with the terminal background time $\chi_f$ free and the corresponding
clock cost minimized. The Pontryagin transversality condition is
$H|_{\chi=\chi_f}=0$ [equation~\eqref{eq:transversality}]---a free-final-clock condition,
\emph{not} a two-event condition. We accordingly reserve ``fixed endpoints'' for a
two-point boundary-value problem (BVP), in which both spatial endpoint positions are
prescribed while the branch-dependent arrival clock is determined by the optimal
solution and remains the quantity being compared.

\paragraph{Standing hypotheses.} All statements labelled I.1--I.6 are made on a
\emph{regular domain}, by which we mean the following three conditions. They are the
precise form of the qualitative requirements listed above, and they are referred to
collectively as (H).

\begin{hypothesis}[timelike selector]\label{hyp:selector}
The selector $W$ is timelike on the open domain $\Omega$ under consideration, so that the
rail constraint $-u\cdot W=\Ehat$ cuts the unit-mass shell in a future-timelike locus.
\end{hypothesis}

\begin{hypothesis}[above freezing]\label{hyp:convex}
The rail energy exceeds the selector norm pointwise on $\Omega$:
$\Ehat>|W|$, where $|W|=\sqrt{-W\cdot W}$.
\end{hypothesis}

\begin{hypothesis}[uniform margin]\label{hyp:coercive}
$\Omega$ is compact and the margin in Hypothesis~\ref{hyp:convex} is uniform,
$\inf_\Omega(\Ehat-|W|)>0$, the metric and the rail one-form being smooth on $\Omega$.
\end{hypothesis}

Hypotheses~\ref{hyp:selector}--\ref{hyp:coercive} are scalar inequalities, directly
checkable on any candidate domain; the geometric properties actually used below---that the
velocity indicatrix is a compact, strictly convex oval, and that the arrival-time
functional is two-sidedly bounded---are \emph{consequences}, derived in
Lemma~\ref{lem:indicatrix}. Hypothesis~\ref{hyp:convex} is what fails at the freezing
surface, Hypothesis~\ref{hyp:coercive} additionally at the horizon, and
Hypothesis~\ref{hyp:selector} inside the ergosphere of the rotating case treated in
Paper~II. The three loci are distinct and are never conflated below.

\begin{auxlemma}[rail indicatrix]\label{lem:indicatrix}
Fix an event at which $W$ is timelike and let $\Ehat>0$. The set
$\mathcal I=\{u:\ u\cdot u=-1,\ u \text{ future-directed},\ -u\cdot W=\Ehat\}$ is
\begin{enumerate}
\item[(i)] empty if $\Ehat<|W|$;
\item[(ii)] the single vector $u=W/|W|$ if $\Ehat=|W|$ \emph{(freezing)};
\item[(iii)] if $\Ehat>|W|$, a round sphere of radius $\sqrt{\Ehat^2/|W|^2-1}$ inside the
spacelike hyperplane $W^\perp$, hence compact, smooth and strictly convex.
\end{enumerate}
In case (iii) the induced set of physical velocities per unit selector time is the image of
that sphere under an invertible affine map, hence a compact, strictly convex ellipse: this
is~\eqref{eq:ellipse}, and its semi-axes vanish exactly as $\Ehat\downarrow|W|$.
Consequently Hypotheses~\ref{hyp:selector}--\ref{hyp:convex} imply the convexity properties
used throughout, and Hypothesis~\ref{hyp:coercive} implies the two-sided
bound~\eqref{eq:coercive} with $\rho,R$ the (uniform) inradius and circumradius of the
resulting ovals.
\end{auxlemma}

\begin{proof}
Set $n=W/|W|$, a unit timelike vector, and split $u=\gamma\,n+\bar w$ with
$\bar w\perp n$ spacelike. The rail constraint reads $-u\cdot W=|W|\,\gamma=\Ehat$, which
fixes $\gamma=\Ehat/|W|$ with no freedom left; the normalization
$u\cdot u=-\gamma^2+|\bar w|^2=-1$ then forces
\[
   |\bar w|^{2}=\gamma^{2}-1=\frac{\Ehat^{2}}{|W|^{2}}-1 .
\]
This is negative, zero or positive according as $\Ehat<|W|$, $\Ehat=|W|$ or $\Ehat>|W|$,
giving (i), (ii) and (iii); in case (iii) the locus is the sphere of radius
$\sqrt{\gamma^2-1}$ in the Euclidean space $W^\perp$, which is compact and strictly convex.
Passing to physical velocities divides the $W^\perp$ components by the fixed component
$\gamma$ and re-expresses them in the coordinate frame; both operations are invertible
linear maps at the given event, and the composition with the translation generated by the
non-orthogonality of $\partial_\chi$ and $n$ is an invertible affine map. The affine image
of a sphere is an ellipsoid, and affine maps preserve compactness and strict convexity, so
the physical velocity locus is a compact strictly convex ellipse whose centre is the
Zermelo wind. Its semi-axes are proportional to $\sqrt{\gamma^2-1}$ and hence vanish as
$\Ehat\downarrow|W|$. Finally the arrival-time Finsler function is the support-type
function of that ellipse, so its inradius and circumradius furnish exactly the constants
$\rho$ and $R$ of~\eqref{eq:coercive}; under Hypothesis~\ref{hyp:coercive} the margin is
uniform on the compact $\Omega$ and so are $\rho>0$ and $R<\infty$.
\end{proof}

\begin{remark}[the Vaidya case, explicitly]\label{rem:vaidya-indicatrix}
For ingoing Vaidya with the Kodama selector $W=K=\partial_v$ one has $|W|=\sqrt f$, so
Hypothesis~\ref{hyp:convex} reads $w\equiv\Ehat^2-f>0$. Substituting
$u=u^v\big(\partial_v+V^r\partial_r+(V^\varphi/r)\partial_\varphi\big)$ into
$g(u,u)=-1$ and $-u_v=\Ehat$ and eliminating $u^v$ gives, with no further assumption,
\[
   \big(V^r+(\Ehat^2-f)\big)^2+\Ehat^2\,(V^\varphi)^2=\Ehat^2\,w ,
\]
an ellipse centred at the ingoing radial wind $f-\Ehat^2<0$ with semi-axes $\Ehat\sqrt w$
and $\sqrt w$: this is precisely the parametrization~\eqref{eq:vaidya-ind} of
section~\ref{sec:vaidya}, and the freezing surface is $w=0$. Because the maximum of a
linear form over an ellipse of centre $c$ and semi-axes $(a_1,a_2)$ is
$p\cdot c+\sqrt{(a_1p_1)^2+(a_2p_2)^2}$, the support function of this oval is exactly the
branch Hamiltonian~\eqref{eq:Hv}, which is Theorem~\ref{thm:rail}(ii) made explicit. The
elimination, the match with~\eqref{eq:vaidya-ind} and the identification of the support
function with~\eqref{eq:Hv} were checked symbolically in \textsc{Mathematica}, all three
residuals reducing to zero identically.
\end{remark}

\Ithm{1}
\begin{theorem}[controlled-rail reduction]\label{thm:rail}
Assume (H). Then minimizing the arrival clock subject to the rail
constraint~\eqref{eq:ellipse} is a time-optimal control problem with control
$\theta\in S^1$, and:
\begin{enumerate}
\item[(i)] for every nonzero costate $p$ the maximizer of the support function
$p\cdot\dot x(\theta)$ over the indicatrix is unique, so the optimal control
$\theta^\star(x,p)$ is single-valued and smooth and no Filippov relaxation is required;
\item[(ii)] Pontryagin's maximum principle delivers the branch Hamiltonian
$H(x,p)=\max_\theta p\cdot\dot x(\theta)$ in closed form;
\item[(iii)] $p_\varphi=J$ is conserved, and the free-arrival-clock transversality
condition forces the branch Hamiltonian to vanish at the free endpoint,
\begin{equation}
    H\big|_{\chi=\chi_f}=0 \qquad\text{(free arrival time).}
    \label{eq:transversality}
\end{equation}
\end{enumerate}
\end{theorem}

\begin{proof}
By Hypothesis~\ref{hyp:convex} the velocity indicatrix~\eqref{eq:ellipse} is a compact,
strictly convex \emph{closed curve} (the boundary of a convex region), and the control is
its direction $\theta\in S^1$, itself a compact one-dimensional set. Maximizing the
support function $p\cdot\dot x(\theta)$ over this curve has a \emph{unique} maximizer by
strict convexity, which is (i): the control set $S^1$ is already compact and the maximizer
single-valued, so \emph{no} Filippov relaxation is needed. This is not merely an
observation about our indicatrix: it is exactly the situation of Filippov's
\cite[\S III]{Filippov1962}, where a \emph{strictly} convex velocity set---one whose
every support plane meets it in a single point---together with a one-to-one
correspondence between velocities and controls is shown to make the maximizer, and
hence the optimal control, depend continuously on the costate. Hypothesis~(H)
supplies both conditions, so continuity of $\theta^\star$ is a published theorem and
only its smoothness is ours. We do \emph{not} claim the
velocity locus is a convex \emph{subset} of the plane, nor convexify it (which would admit
averaged velocities violating the unit-mass shell). Excluding $p=0$---legitimate by the
normality statement of Theorem~\ref{thm:existence}, which is proved independently
below---the optimal control $\theta^\star(x,p)$ is single-valued and smooth away from the
freezing and horizon degeneracies, where Hypothesis~\ref{hyp:coercive} fails and which are
treated separately. Pontryagin's maximum principle (PMP) \cite{Pontryagin,Liberzon} then
delivers the branch Hamiltonian $H(x,p)=\max_\theta p\cdot\dot x(\theta)$, evaluated in
closed form in section~\ref{sec:vaidya}, which is (ii). Finally $\varphi$ is cyclic, so
$p_\varphi=J$ is conserved, and the terminal cost being the free arrival clock the
transversality condition of the PMP reads $H|_{\chi=\chi_f}=0$, which
is~\eqref{eq:transversality} and completes (iii).
\end{proof}

\begin{remark}[vakonomic, not nonholonomic]
This is a genuine time-optimal control problem, for which the Pontryagin
(vakonomic-type) reduction is the correct prescription; we do not invoke the nonholonomic
d'Alembert principle, whose equations differ from the vakonomic ones unless the constraint
distribution is integrable---a coincidence that convexity does not entail.
\end{remark}

\begin{remark}[the branch Hamiltonian is not a first integral]
For a non-autonomous system $H$ is \emph{not} a global first integral: along the
extremal $\dd H/\dd\chi=\partial_\chi H$, proportional to $m'(v)$ (Vaidya) or $A'(\eta)$
(Thakurta--Kerr), so $H\neq0$ in the interior---consistently, $p_v(r_1)=0$ at the
free arrival radius yet $p_v(r)\neq0$ along the path (figure~\ref{fig:vaidya-traj}).
The globally conserved object is the extended Hamiltonian $\mathcal H=p_\chi+H=0$ on
the enlarged phase space, with $p_\chi$ conjugate to the background coordinate $\chi$; only
in the stationary limit is $H$ autonomous and $H\equiv0$ a genuine first integral.
\end{remark}

\paragraph{Cost of control.} For an accelerated worldline with rail four-force
$a^\mu=Du^\mu/\dd\tau$, differentiating the invariant gives the \emph{exact}
identity
\begin{equation}
    \frac{\dd(-u\cdot W)}{\dd\tau}=-\,a\cdot W-u^a u^b\,\nabla_{(a}W_{b)},
    \label{eq:cost}
\end{equation}
whose two terms play distinct roles. The constraint $a\cdot u=0$ holds
\emph{automatically} for any four-force that preserves the mass shell
$u\cdot u=-1$ (it is the derivative of the shell condition); it is \emph{not} a
special ``magnetic'' property and does not by itself imply that the rail performs no
work relative to a physical observer. It means only that no rest mass is created;
the rail is nonetheless free to exert power \emph{against the selector} $W$ (the term
$a\cdot W$). One should distinguish (i) preservation of rest mass ($a\cdot u=0$),
(ii) power measured by a chosen observer congruence, (iii) the constraint force
itself, and (iv) the additional power needed to compensate non-stationarity.
A free geodesic ($a=0$) would let $-u\cdot W$ drift at $-u^au^b\nabla_{(a}W_{b)}$,
which vanishes for all $u$ iff $W$ satisfies the Killing equation; to hold
$-u\cdot W=\Ehat$ fixed the rail supplies exactly $a\cdot W=-u^au^b\nabla_{(a}W_{b)}$,
so that the left-hand side of~\eqref{eq:cost} is zero. This control power equals the
failure of the Killing equation---the missing symmetry---which is why the
non-stationary rail is legitimate but not free: $\Ehat$ is actively controlled, not
conserved by a symmetry. The optical reductions of $H$
reproduce the Randers/Zermelo structure
\cite{GibbonsRanders2009,CaponioWind}, and the sense of ``reproduce'' should be
stated exactly. Gibbons, Herdeiro, Warnick and Werner establish a triality between
Zermelo navigation, Randers--Finsler geodesics and the \emph{null} geodesic flow of
a conformally stationary metric, with Randers data $a_{ij}=V^{-2}\gamma_{ij}$ and
$b_i=-\omega_i$ read off the threading form
\cite[Eqs.~(19)--(23)]{GibbonsRanders2009}; their $a_{ij}$ is the Fermat optical
metric. The controlled rail reproduces that \emph{form}---a Riemannian term plus a
one-form linear in the velocity---but with Perlick's fixed-energy factor in place of
the null one, so the coefficients agree only in the limit $\Ehat\to\infty$. What is
inherited unchanged is the structural consequence recorded in
Paper~II: the Randers data of a conformally stationary metric are conformal
invariants, and the Finsler condition $|b|_a<1$ that makes $F$ a genuine norm fails
exactly on the ergosurface.

Caponio, Javaloyes and S\'anchez enter at exactly that failure, and the
distinction is worth drawing because it marks where our problem stops and a
developed theory continues. Their \emph{wind Finslerian structures} generalise
Finsler and Kropina metrics to indicatrices that need not contain the zero vector,
which models Zermelo navigation ``even when the trajectories of the moving objects
\dots\ are influenced by \emph{strong} winds'', and simultaneously describes the
causal structure of spacetimes carrying a non-vanishing Killing field
\cite{CaponioWind}. Their treatment splits along precisely the boundary that
governs our control domain: causal $K$ gives Randers--Kropina metrics
\cite[\S4]{CaponioWind}, arbitrary $K$ requires the general wind Riemannian
structures \cite[\S5]{CaponioWind}, and the surface separating them appears there
as a $K$-horizon \cite[\S8.3]{CaponioWind}. So the geometry does not end at
$r=r_e$; what ends is the Randers description and, with it, the compact-control
problem of Paper~II. The continuation is a wind Riemannian structure, and we do not
pursue it here.

\paragraph{Operational reading: thrust.} Maintaining the controlled invariant generally
requires nonzero \emph{proper} acceleration $a^\mu=Du^\mu/\dd\tau$ ($a=0$ only on a
free-fall geodesic), so the rail force is the thrust per unit mass, spacelike
($a\cdot u=0$) with $|a|$ diverging as the freezing/horizon surface is approached. The
present variational problem minimizes the selected elapsed \emph{clock}, not thrust, fuel
or integrated control effort, which define distinct optimal-control problems. In the
non-stationary case the $a\cdot W$ term is the proper acceleration that maintains the rail
against the non-stationary mass-rate drift.

\paragraph{Geometric dictionary.} The indicatrix~\eqref{eq:ellipse} carries two
independent data---its \emph{center} $c=(c_r,c_\varphi)$ (the Zermelo wind) and
the \emph{eccentricity} of its semi-axes (the anisotropy of the local speed,
$a_r\neq a_\varphi$)---and each governs a distinct feature of the
brachistochrones, the same for every branch:
\begin{itemize}\itemsep2pt
\item eccentricity $\neq0$ (a spatial gradient $\partial_r f\neq0$)
$\Rightarrow$ the $t$ and $\tau$ branches split in depth, $n_t/n_\tau=E/f\neq1$;
\item azimuthal offset $c_\varphi\neq0$ (frame dragging)
$\Rightarrow$ rotational plunge inversion becomes possible;
\item radial offset $c_r\neq0$ (a non-stationary wind, $m'\neq0$)
$\Rightarrow$ endpoint-sensitive memory, $p_v\neq0$ along the path;
\item shrinking axes $R\to0$ (local speed $v\to0$) $\Rightarrow$ freezing and
capture.
\end{itemize}
Eccentricity and center offset are \emph{separate} levers: the first switches on
the $t/\tau$ split, the second the inversions. In particular the three branches
\emph{degenerate}---the $t$, $\tau$ and $\eta$ brachistochrones coincide---iff the
indicatrix is, up to a position-independent conformal scale, a rigid centered
circle ($c=0$, $a_r=a_\varphi$, radius depending on the flow parameter alone),
equivalently iff the optical cost ratios $n_t/n_\tau,\,n_t/n_\eta$ are
position-independent. This is exactly spatial homogeneity: FLRW is the unique
such case in our ladder (section~\ref{sec:flrw}), whereas any gradient or wind
breaks it---even the static, non-rotating Schwarzschild ellipse, which is
centered ($c=0$) yet eccentric, already splits $t$ from $\tau$. The indicatrix
shape is thus the organizing thread of the paper: FLRW (circle) $\to$ Vaidya
(radial wind) $\to$ Thakurta--Kerr (angular wind) is a progression of ellipse
deformations, each switching on one further phenomenon.

\begin{figure}[tbp]
    \centering
    % ALT: Two panels of closed curves in the local velocity plane: left, nested circles
    %      centred on the origin and shrinking; right, ellipses displaced toward negative
    %      radial velocity and elongated radially.
    \includegraphics[width=\textwidth]{Immagini/fig_indicatrici.pdf}
    \caption{Rail indicatrices (linear constraint $\cap$ normalization
    $=$ ellipse) in the two spherically symmetric applications of this paper.
    Left: spatially flat FLRW gives centered shrinking circles, so the three
    clock branches remain degenerate. Right: ingoing Vaidya produces a radial
    wind and an offset, eccentric indicatrix, switching on branch splitting and
    endpoint-sensitive memory.}
    \label{fig:indicatrici}
\end{figure}

\subsection{Existence, normality, and global optimality}
\label{sec:control-foundations}
The Pontryagin reduction above supplies only \emph{necessary} conditions. We
record here the three complementary facts that make the constrained problem well
posed for the non-stationary rail: an existence theorem for the minimizer, the
normality of every regular extremal (so the reduction $-u\cdot W=\Ehat$ is never
vacuous), and a Hamilton--Jacobi certificate that upgrades a saturated extremal to
a \emph{global} optimum. All three are elementary under the standing hypothesis
that the indicatrix is a compact, uniformly convex oval [equation~\eqref{eq:ellipse}].

\Ithm{2}
\begin{theorem}[existence and normality; direct method]\label{thm:existence}
Assume (H), and let the spatial arc range over the compact regular domain $\Omega$. By
Lemma~\ref{lem:indicatrix} the arrival-time Finsler function then obeys the uniform
two-sided bound
\begin{equation}
  \rho\,|\nu| \le F(x,\nu) \le R\,|\nu|, \qquad 0<\rho\le R<\infty,
  \label{eq:coercive}
\end{equation}
with $\rho,R$ the inradius and circumradius of the ovals. Then:
\begin{enumerate}
\item[(i)] \emph{(existence)} the arrival time
$T[\gamma]=\int F(\gamma,\dot\gamma)\,\dd\sigma$ attains its infimum in the class of
Lipschitz arcs joining the fixed endpoints;
\item[(ii)] \emph{(normality)} every regular minimizer is a \emph{normal} extremal: the
abnormal multiplier $p_0$ cannot vanish, hence normalizes to $-1$, and the costate $p$
never vanishes along the extremal.
\end{enumerate}
\end{theorem}

\begin{proof}
(i) Parametrize spatial arcs on $[0,1]$. The clock along an arc is not a datum but the
solution of an ordinary differential equation: by construction
\begin{equation}
   \frac{\dd\chi}{\dd\sigma}=F\big(x(\sigma),\dot x(\sigma),\chi(\sigma)\big),
   \qquad \chi(0)=\chi_0 ,
   \label{eq:clock-ode}
\end{equation}
and the cost is $T[x]=\chi(1)-\chi_0$. We argue in three steps.

\emph{A priori bounds.} Let $T_0$ be the cost of any admissible competitor. Along a
minimizing sequence, $\rho\,\mathrm{length}(x_n)\le T[x_n]\le T_0$ by the left inequality
of~\eqref{eq:coercive}, so the lengths are bounded by $\Lambda=T_0/\rho$. Reparametrizing
each $x_n$ proportionally to arclength gives $|\dot x_n|\equiv\mathrm{length}(x_n)\le\Lambda$
a.e., a uniform Lipschitz bound. The right inequality gives
$0\le\dd\chi/\dd\sigma\le R\Lambda$, so the clocks $\chi_n$ are uniformly Lipschitz and
uniformly bounded, and remain in a compact interval $[\chi_0,\chi_0+R\Lambda]$ on which,
with $\Omega$ compact, $F$ is smooth.

\emph{Gr\"onwall stability of the clock.} On $\Omega\times[\chi_0,\chi_0+R\Lambda]$ the
function $F$ is $C^1$, hence Lipschitz in its slow argument: there is $L<\infty$ with
$|F(x,\nu,\chi)-F(x,\nu,\chi')|\le L\,|\nu|\,|\chi-\chi'|$, the factor $|\nu|$ appearing
because $F$ is positively homogeneous of degree one in $\nu$. Let $\chi,\tilde\chi$ solve
\eqref{eq:clock-ode} for the same arc but different initial data, or for two arcs with the
same velocity field; subtracting and using the above,
\[
   \big|\chi(\sigma)-\tilde\chi(\sigma)\big|
   \le\big|\chi(0)-\tilde\chi(0)\big|
   +L\Lambda\!\int_0^{\sigma}\!\big|\chi-\tilde\chi\big|\,\dd\sigma' ,
\]
so Gr\"onwall's inequality gives
$\|\chi-\tilde\chi\|_\infty\le e^{L\Lambda}|\chi(0)-\tilde\chi(0)|$. In particular
\eqref{eq:clock-ode} has a unique solution for each arc, the solution map
$x\mapsto\chi_x$ is well defined, no trajectory escapes the compact set in which the
bounds were derived, and $T$ depends on the arc only through that map.

\emph{Existence.} Pass to the augmented state $y=(\chi,x)$ and the extended velocity set
\[
   \mathcal V(y)=\big\{(t,\nu):\ |\nu|\le\Lambda,\ \ F(x,\nu,\chi)\le t\le R\Lambda\big\},
\]
which is nonempty, compact and---because $F$ is convex in $\nu$, so that its epigraph is
convex---\emph{convex}, and depends continuously on $y$ by the smoothness just used. The
relaxed problem of minimizing $\chi(1)$ over trajectories of the differential inclusion
$\dot y\in\mathcal V(y)$ therefore satisfies the hypotheses of the Filippov--Cesari
existence theorem \cite{Cesari1983,Filippov1962}, and admits a minimizer. It is worth
recording that each hypothesis of \cite[Thm.~1]{Filippov1962} is met with margin. The
velocity set is required only to be \emph{upper semicontinuous with respect to
inclusion}, whereas ours depends continuously on $y$. The growth condition
$y\cdot\dot y\le C(|y|^2+1)$ is immediate from the uniform bounds
$|\nu|\le\Lambda$ and $0\le\dd\chi/\dd\sigma\le R\Lambda$ established above, which make
$\mathcal V$ uniformly bounded outright. Nonemptiness of the admissible class---an
explicit hypothesis there, not a conclusion---holds because $\Omega$ is a connected
regular domain, so its fixed endpoints are joined by at least one Lipschitz arc.
Finally, the state constraint $x\in\Omega$ places the problem in Filippov's general
formulation \cite[\S IV]{Filippov1962}, which requires the admissible region to be
\emph{closed} precisely so that the limit of a minimizing sequence of curves lying in
it again lies in it; the compactness of $\Omega$ assumed in the theorem is that
requirement, and is the reason the domain is not taken to be the open exterior. Since $\chi(1)$
is monotone in $t$, any minimizer saturates $t=F(x,\nu,\chi)$ almost everywhere, hence
solves~\eqref{eq:clock-ode} and is an admissible arc of the original problem with the same
cost. Coercivity forbids the infimum from being attained by escape, and the two problems
have the same infimum, so the minimizer of the relaxed problem is a minimizer of $T$.

(ii) Suppose $p_0=0$. The maximized Hamiltonian would then read $\max_\nu p\cdot\nu=0$
with $p\neq0$; but the support function of the oval obeys
$h(p)=\max_{\nu\in\partial B}p\cdot\nu\ge\rho\,|p|>0$ by the left inequality
of~\eqref{eq:coercive}, a contradiction. Hence $p_0$ normalizes to $-1$ and $p$ never
vanishes.
\end{proof}

\begin{remark}[relaxation of the epigraph is not convexification of the indicatrix]
\label{rem:no-convexification}
The convexity invoked in the existence proof is that of the \emph{epigraph}
$\{t\ge F(x,\nu,\chi)\}$ in the augmented variable $t$---the classical Cesari property---and
it is used only to close the minimizing sequence. It is \emph{not} a convexification of the
velocity indicatrix, which Theorem~\ref{thm:rail} deliberately avoids: replacing the oval
by its convex hull would admit averaged velocities violating the unit-mass shell, whereas
enlarging $t$ upwards only makes a competitor slower and is therefore discarded by the
minimization itself (the saturation step above). The two statements are consistent: no
relaxation survives in the minimizer.

That the epigraph convexity cannot simply be dropped is shown by Filippov's own
counterexample \cite[\S V]{Filippov1962}: for
$\dot x=-y^2+u^2,\ \dot y=u,\ |u|\le1$ steered from the origin to $(1,0)$ in minimal
time, the velocity set is not convex and \emph{no} optimal control exists---every
minimizing sequence chatters between $u\simeq+1$ and $u\simeq-1$ and converges to a
sliding regime whose velocity lies in the convex hull of the velocity set but not in
the set itself. The distinction drawn in this remark is therefore the operative one:
convexifying the \emph{indicatrix} would manufacture exactly such sliding
velocities---superluminal averages violating the unit-mass shell---while convexifying
the \emph{epigraph} in the augmented clock variable adds only slower competitors and
is undone by the saturation step.
\end{remark}

\begin{remark}[what normality buys, and where the hypotheses fail]
Part (ii) is what makes the reduction to the branch Hamiltonian
$H=\max_\theta p\cdot\dot x(\theta)$ together with the rail invariant $-u\cdot W=\Ehat$
legitimate along the \emph{entire} extremal, and it is the control-theoretic content of
``excluding $p=0$'' in the proof of Theorem~\ref{thm:rail}: nontriviality is not an
assumption but a consequence of $\rho>0$. The two-sided bound~\eqref{eq:coercive} fails
exactly at the freezing/horizon surfaces, where $\rho\to0$ and the velocity indicatrix of
section~\ref{sec:indicatrix} degenerates; Theorem~\ref{thm:existence} therefore applies on
any compact subdomain bounded away from them.
\end{remark}

\Ithm{3}
\begin{proposition}[Hamilton--Jacobi verification]\label{prop:hjb}
Assume (H). The value function of the non-autonomous problem obeys, in the cost-to-go
convention, the Hamilton--Jacobi--Bellman equation
\begin{equation}
  \partial_\chi V+\min_{\theta\in S^1}\big\{\,\ell_{\mathcal T}(x,\theta,\chi)
  +\nabla_x V\cdot f(x,\theta,\chi)\,\big\}=0
  \label{eq:hjb}
\end{equation}
(equivalently, an autonomous HJB on the augmented state $y=(\chi,x)$; choosing
$\mathcal T=\chi$ gives $\ell_{\mathcal T}=1$ and the usual minimum-time form).
Let now $V$ be a Lipschitz \emph{calibration} on a region where it is single-valued,
that is, a viscosity subsolution of~\eqref{eq:hjb} in the forward (accumulated-cost)
convention:
\begin{equation}
  \partial_\chi V+\nabla_x V\cdot f(x,\theta,\chi)\;\le\;\ell_{\mathcal T}(x,\theta,\chi)
  \qquad\text{for \emph{every} }\theta\in S^1 .
  \label{eq:calibration}
\end{equation}
If a candidate PMP extremal \emph{saturates}~\eqref{eq:calibration}---equality along it,
which in the stationary case is $\nabla V\cdot\dot x=F$ pointwise---then that extremal
realizes the value function and is a \emph{global} minimizer on that region.
\end{proposition}

\begin{proof}
The maximum principle is necessary but not sufficient: a PMP extremal is a
\emph{candidate} minimizer only. Minimality implies the PMP, but the PMP implies
neither global nor \emph{local} minimality---that would need a second-variation
or conjugate-point argument (Theorem~\ref{thm:jacobi}) or another sufficient
condition. The sufficient global certificate is the
calibration~\eqref{eq:calibration}. Let $(x(\cdot),\theta(\cdot))$ be
\emph{any} admissible competitor joining the data. Since~\eqref{eq:calibration} holds for
every $\theta$, it holds in particular for the competitor's own control, so
\[
   \frac{\dd}{\dd\chi}\,V\big(\chi,x(\chi)\big)
   =\partial_\chi V+\nabla_x V\cdot f\big(x,\theta(\chi),\chi\big)
   \;\le\;\ell_{\mathcal T}\big(x,\theta(\chi),\chi\big).
\]
Integrating from $\chi_0$ to the arrival $\chi_f$ gives
$V(\chi_f,x_f)-V(\chi_0,x_0)\le\int\ell_{\mathcal T}\,\dd\chi=J_{\mathcal T}$: the
increment of $V$ between the prescribed data is a lower bound for the cost of every
competitor. Along the candidate the inequality is an equality by hypothesis, so its cost
equals that lower bound and no competitor can beat it. Finally
$V(\chi_f,x_f)-V(\chi_0,x_0)$ therefore equals the value function, and the candidate
realizes it. In the stationary case~\eqref{eq:calibration} reduces to
$F^\ast(x,\nabla V)\le1$ with $F^\ast$ the dual (optical) norm, the familiar statement
that $V$ is a subsolution of the eikonal equation, and saturation is
$\nabla V\cdot\dot x=F$.
\end{proof}

\begin{remark}[scope of the certificate]
We do \emph{not} construct such a global $V$ for generic non-stationary Kerr/Vaidya
trajectories; in the stationary limit the closed-form optical (Randers) metric of
section~\ref{sec:vaidya} supplies one on any region carrying a single-valued $V$, and
elsewhere our statements are those of PMP extremals with local checks. Where $V$ turns
multivalued the certificate localizes, as follows.
\end{remark}

\paragraph{Conjugate points, Maxwell set, cut locus.} Two mechanisms end
optimality. A \emph{conjugate point} (Jacobi) ends \emph{local} minimality when
the second variation degenerates; up to the first conjugate point the extremal is
a strict local minimum. A \emph{Maxwell point}---two distinct extremals of equal
arrival time---ends \emph{global} minimality earlier and is exactly a shock
(gradient discontinuity) of $V$, i.e.\ a point of the cut locus. Strict convexity
of the indicatrix removes Filippov relaxation but does \emph{not} exclude Maxwell
points: distinct windings past the rotation axis, or the two $\pm J$ arcs reaching
a common target, can tie. A closed-form extremal admitting a single-valued saturated
verification function is globally certified only up to its first Maxwell or conjugate
obstruction---the first such singularity of $V$. For the generic non-stationary
trajectories treated here we establish PMP extremality and local numerical minimality
checks, \emph{not} a complete global Hamilton--Jacobi certificate; beyond the first
obstruction the global optimum is selected by comparing arrival times across the competing
extremals. On the frozen family both obstructions are located exactly in
section~\ref{sec:vaidya-jacobi}: no conjugate point occurs before the turning point, so a
Mayer field exists on every monotone arc and the criterion above does apply there, whereas
an arc through periapsis develops a genuine caustic at a computable rail energy, and the
Maxwell set is the locus of axis crossings.

\paragraph{What is \emph{not} claimed inside the ergosphere.} Where the Killing
selector $\xi$ becomes spacelike---inside the ergosphere---the invariant
$-u\cdot\xi=\Ehat$ no longer bounds a timelike-future cost and the indicatrix oval
degenerates: a closed-form $\varphi(r)$ continued through $\re$ is a genuine
\emph{analytic continuation} of the
extremal, not a certified physical minimum. The existence and Hamilton--Jacobi
statements above are asserted only outside the ergosphere, on the region where $F$
is a bona fide convex indicatrix; inside, a timelike observer congruence or a different
timelike selector would be required to formulate a compact future-timelike indicatrix,
which we do not develop. Freezing (where
$\rho\to0$), turning (where a semi-axis vanishes), and the separatrix (a double
root of the radical) are distinct degeneracies and are not conflated.

\subsection{Selecting $W$: the Killing / conformal-Killing / Kodama hierarchy}
\label{sec:whierarchy}
The selector $W$ is fixed by a hierarchy of decreasing symmetry, each level
recovering the previous one in the appropriate limit.

\emph{(a) Killing.} If a timelike Killing vector $\xi$ exists, $W=\xi$ and the
right-hand side of equation~\eqref{eq:cost} vanishes identically: $-u\cdot\xi$ is
Noether-conserved and the rail is free. This is Perlick's stationary case
\cite{Perlick1991}.

\paragraph{The time-dependent optical metric, and Perlick's as its stationary member.}
The rail does more than supply extremal equations: it produces an \emph{optical metric},
and it does so whether or not the geometry is stationary. This is the precise sense in
which the construction extends Perlick's.

\begin{auxtheorem}[time-dependent optical metric]\label{thm:optical}
Assume (H1)--(H2), let $\chi$ be adapted to the selector ($W=\partial_\chi$), write
$f=|W|^2=-g_{\chi\chi}$, and put the metric in threading form
$g=-f(\dd\chi-\omega_a\dd x^a)^2+h_{ab}\dd x^a\dd x^b$, where $f,\omega_a,h_{ab}$ may all
depend on $\chi$. Then the rail arrival-time functional is the Randers-type Finsler
function
\begin{equation}
   F(\chi,x,\nu)=\beta_a\nu^a+\sqrt{a_{ab}\nu^a\nu^b},
   \qquad \beta_a=\omega_a,\qquad
   a_{ab}=\frac{\Ehat^2\,h_{ab}}{f\,(\Ehat^2-f)},
   \label{eq:optical-metric}
\end{equation}
and
\begin{enumerate}
\item[(i)] $a_{ab}$ is positive definite \emph{if and only if} $\Ehat>|W|$: Hypothesis
\ref{hyp:convex} \emph{is} the positive-definiteness of the optical metric, whose conformal
factor $\Ehat^2/[f(\Ehat^2-f)]$ diverges as $\Ehat\downarrow|W|$, the freezing surface;
\item[(ii)] $\|\beta\|_a^2=\big[f(\Ehat^2-f)/\Ehat^2\big]\,|\omega|_h^2$, strictly
increasing in $\Ehat$ with limit $f|\omega|_h^2$; hence the Randers condition
$\|\beta\|_a<1$ is \emph{implied by} the null Fermat--Randers condition $f|\omega|_h^2<1$,
for every $\Ehat$;
\item[(iii)] if the data are $\chi$-independent, $F$ is Perlick's fixed-energy optical
metric \cite{Perlick1991}, and $\Ehat\to\infty$ gives the Fermat--Randers metric
$h_{ab}/f$ \cite{GibbonsWerner2008}.
\end{enumerate}
\end{auxtheorem}

\begin{proof}
Solving the rail $-u\cdot W=\Ehat$ jointly with $g(u,u)=-1$ for $\chi'=\dd\chi/\dd\sigma$
gives a quadratic; substituting the right-hand side of~\eqref{eq:optical-metric} into it
annihilates it, the rational part and the coefficient of the radical vanishing separately,
and $f\,F-\omega_a\nu^a f=f\sqrt{a_{ab}\nu^a\nu^b}>0$ identifies it as the
future-directed root. The computation is \emph{pointwise in the event}: no derivative of
the metric appears, so the $\chi$-dependence of $(f,\omega,h)$ never enters, and the same
algebra that yields Perlick's formula yields~\eqref{eq:optical-metric} verbatim. Part (i)
is immediate since $h_{ab}$ is positive definite; part (ii) follows from
$a^{ab}=f(\Ehat^2-f)h^{ab}/\Ehat^2$ and
$\dd\big[f(\Ehat^2-f)/\Ehat^2\big]/\dd\Ehat=2f^2/\Ehat^3>0$; part (iii) is the
specialization.
\end{proof}

\begin{remark}
Two consequences are worth separating. First, the extension is not an analogy: the
\emph{same} closed formula covers both regimes, and stationarity enters only by removing
the $\chi$-argument. Second, it is not a reduction to Finsler geodesics either---because
$F$ depends on $\chi$, the brachistochrone is \emph{not} a geodesic of a fixed Randers
metric, which is exactly the non-autonomy that forces the Hamilton--Jacobi criterion of
Proposition~\ref{prop:hjb} in place of Finsler geodesic theory. Perlick's case is the one
where that argument collapses to the standard theory. Part (ii) adds that the rail
structure is no worse behaved than the null one: wherever Fermat--Randers optics is
defined, so is the rail optics, at every rail energy.
\end{remark}

\paragraph{Recovery of Perlick's stationary optics, explicitly.} In the Killing limit
$W=\xi=\partial_t$ the controlled-rail arrival-time Finsler function is exactly a
Randers metric, so the construction reduces to Perlick's fixed-energy optics rather
than merely resembling it. Solving the rail $-u\cdot\xi=\Ehat$ jointly with
$g(u,u)=-1$ for the coordinate-time rate $t'=\dd t/\dd\sigma$ gives a quadratic whose
future-directed root is
\begin{equation}
F_{\rm rail}=\beta_a \nu^a+\sqrt{a_{ab}\,\nu^a \nu^b},\quad
\beta_a=\frac{g_{ta}}{f},\quad
a_{ab}=\frac{\Ehat^2\bigl(g_{ta}g_{tb}+f\,g_{ab}\bigr)}{f^2(\Ehat^2-f)},
\label{eq:perlick-randers}
\end{equation}
with $f=-g_{tt}$ and $\nu^a=\dd x^a/\dd\sigma$. In the threading form
$g=-f(\dd t-\omega_a\dd x^a)^2+h_{ab}\dd x^a\dd x^b$ one has
$g_{ta}g_{tb}+f\,g_{ab}=f\,h_{ab}$, so $\beta_a=\omega_a$ and
$a_{ab}=\Ehat^2 h_{ab}/[f(\Ehat^2-f)]$---precisely Perlick's fixed-energy stationary
optical metric \cite{Perlick1991}. The null/high-energy limit $\Ehat\to\infty$ gives
$a_{ab}\to h_{ab}/f$, the standard Fermat--Randers optical metric of Kerr
\cite{GibbonsWerner2008}. This is the functional equality $F_{\rm rail}=F_{\rm Perlick}$
(stronger than coincidence of the Hamiltonian flows on $H=0$), verified against the
directly solved $F_{\rm rail}$ on equatorial Kerr to machine precision.

\emph{(b) Conformal Killing.} If the spacetime is conformally stationary,
$g=\Omega^2\hat g$ with $\hat g$ stationary, the timelike conformal Killing
vector $\partial_\eta$ selects the conformal-time branch; $-u\cdot W$ is exactly
conserved for null rays and controlled up to the conformal factor for massive
rails. FLRW and constant-$A$ Thakurta--Kerr live at this level.

\emph{(c) Kodama.} In a spherically symmetric \emph{dynamical} spacetime with no
Killing or conformal Killing vector, the Kodama vector
$K^a=\varepsilon^{ab}\nabla_b r$ \cite{Kodama1980,Hayward1996} is the canonical
replacement: it is divergence-free (equation~\eqref{eq:kodama}), is \emph{parallel}
to $\xi$ in the static limit, and its flux $G^{ab}K_b$ is conserved with charge the
Misner--Sharp mass. The qualification matters and is Kodama's own, emphasised by
Abreu and Visser: the Kodama vector ``does \emph{not} in general reduce to the
Killing vector in a static spacetime; all that one can say in general is that in
static spacetimes it is \emph{parallel} to the Killing vector''
\cite[\S III]{AbreuVisser2010}. For a static metric
$-E(r)^2\dd t^2+L(r)^2\dd r^2+r^2\dd\Omega^2$ one finds
\begin{equation}
    K=\frac{1}{E(r)L(r)}\,\partial_t ,
    \label{eq:kodama-norm}
\end{equation}
so $K=\xi$ exactly when $g_{tt}g_{rr}=-1$. Schwarzschild and Vaidya satisfy that,
which is why the identification is safe throughout this paper; a general static
interior, where $e^{\Phi+\Lambda}\neq1$, does not, and there the two differ by that
factor. Equation~\eqref{eq:kodama-norm} is verified in
\texttt{../paper2/verification/verify\_vaidya\_tube.wls}. Vaidya is the paradigm (section~\ref{sec:kodama-energy}).

\emph{(d) Convention.} Absent even spherical symmetry the Kodama construction
fails and $W$ becomes a declared convention, with no distinguished normalization
of the cost one-form. Levels (a)--(c) are what make the non-stationary rail
canonical rather than arbitrary.

\subsection{The price of the rail, and where the optical metric ceases to descend}
\label{sec:price}
The hierarchy above is qualitative. Two quantitative statements make it sharp,
and both are consequences of a single identity.

\begin{auxlemma}[rail drift and the control force]\label{lem:drift}
Along any worldline with four-velocity $u$ and four-acceleration
$a=\nabla_uu$,
\begin{equation}
    \frac{\dd}{\dd\tau}\big(\!-\!g(u,W)\big)
      =-g(a,W)-\tfrac12(\mathcal L_Wg)_{\mu\nu}u^\mu u^\nu .
    \label{eq:drift-general}
\end{equation}
In particular the rail charge is conserved along \emph{geodesics} if and only if
$\mathcal L_Wg$ annihilates $u\otimes u$; and holding the rail on a non-geodesic
forces $g(a,W)=-\tfrac12(\mathcal L_Wg)_{\mu\nu}u^\mu u^\nu$.
\end{auxlemma}

At level (a) the right-hand side vanishes identically and the rail is free, which
is the Noether statement already made. Below level (a) it does not, and
\eqref{eq:drift-general} quantifies the cost. For ingoing Vaidya with the Kodama
selector the only nonvanishing component is $(\mathcal L_Kg)_{vv}=2m'(v)/r$, with
$\nabla\!\cdot\!K=0$, so $K$ is Killing exactly when $m'=0$ and is \emph{not} even
conformal Killing otherwise; along a geodesic
\begin{equation}
    \frac{\dd\Ehat}{\dd\tau}\bigg|_{\rm geodesic}=-\frac{m'(v)\,(u^v)^2}{r} .
    \label{eq:vaidya-drift}
\end{equation}
At level (b), where $\mathcal L_Wg=2\psi g$, the drift is $\psi$ itself,
independent of the four-velocity.

Equation~\eqref{eq:drift-general} also bounds the thrust. Since $g(a,u)=0$, only
the part of $W$ orthogonal to $u$ contributes, and
$|W_\perp|^2=g(W,W)+g(W,u)^2=\Ehat^2-|W|^2$; Cauchy--Schwarz in the rest space of
$u$ gives the \emph{sharp} bound
\begin{equation}
    |a|\ \ge\ \frac{\big|\tfrac12(\mathcal L_Wg)_{\mu\nu}u^\mu u^\nu\big|}
                   {\sqrt{\Ehat^2-|W|^2}},
    \label{eq:thrust-general}
\end{equation}
attained when $a$ is parallel to $W_\perp$---the minimum-thrust control. The
denominator is exactly the quantity that vanishes on the freezing surface
$\Ehat=|W|$ of Hypothesis~\ref{hyp:convex}, so the minimum thrust \emph{diverges}
there, as an inverse square root in the margin $\Ehat^2-|W|^2$. It does not
diverge merely because $|W|\to0$. The two boundaries of the control domain are
therefore of different kinds: one is a loss of compactness at finite cost, the
other a divergence of cost, and only the second makes the rail physically
unattainable before the problem ceases to be posed.

\paragraph{Why Theorem~\ref{thm:optical} must carry a time argument.} The optical
metric of Theorem~\ref{thm:optical} is not merely conveniently written with the
adapted time in it; there is an obstruction to removing it. Projecting along the
orbits of $W$ is a submersion, and giving the base the fixed-energy metric
$a_{ab}=\Ehat^2h_{ab}/[f(\Ehat^2-f)]$ makes it \emph{horizontally conformal} with
dilation
\begin{equation}
    \Lambda^2=\frac{\Ehat^2}{f\,(\Ehat^2-f)},
    \label{eq:dilation-p1}
\end{equation}
the Perlick factor itself. That the rescaling is \emph{conformal} at all is not
an accident: by the Fuglede--Ishihara dichotomy
\cite{Fuglede1978,Ishihara1979,BairdWood2003} a semiconformal map is either
identically zero or conformal on the complement of the kernel of its
differential.

What fails away from level (a) is that the structure should \emph{descend}. A
conformal submersion admits a rescaling of the base turning it into a Clairaut
Riemannian submersion only when its dilation is constant along the fibres
\cite{MeenaZawadzki2024}, and by Bishop's theorem \cite{Bishop1972} a Riemannian
submersion with connected fibres is Clairaut---so that the transverse invariant is
conserved along geodesics---if and only if its fibres are totally umbilical with
mean curvature of gradient form. When $W$ is not Killing,
$\Lambda$ varies along the fibres, the horizontal distribution ceases to be
totally geodesic, and no metric on a quotient reproduces the arrival time. The
optical geometry then exists only on the total space, with the adapted time as an
explicit argument. That is the content of Theorem~\ref{thm:optical}, and the
reason it cannot be stated on a quotient the way Perlick's can.

\begin{remark}
Two structures from the stationary theory are, for the same reason, unavailable
here and are not assumed. The fixed-energy variational theory of
\cite{CaponioCoronaGiamboPiccione2024} makes the set of curves of constant
Noether charge a $C^1$ closed submanifold of the path space---but the argument
uses conservation, which is what fails below level (a). And the Palais--Smale
property established for the constraint $\langle\dot z,\dot z\rangle=a<0$ in
\cite{GiannoniMasielloPiccione1998} is for a constraint the flow preserves, which
the rail is not. Theorem~\ref{thm:existence} accordingly proceeds by the direct
method on the augmented state rather than by restricting a functional to a
constraint submanifold.
\end{remark}


% ======================================================================
\section{FLRW: the degenerate base case}
\label{sec:flrw}
% ----------------------------------------------------------------------
For spatially flat FLRW, $g=a(\eta)^2\eta_{\rm Mink}$, the rail uses the
conformal Killing vector $\partial_\eta$; the local speed is
\begin{equation}
    v(\eta)=\sqrt{1-a(\eta)^2/\Ehat^{2}},
    \label{eq:flrwspeed}
\end{equation}
vanishing at the freezing barrier $a=\Ehat$.
\paragraph{Branch degeneracy.} Because $a=a(\eta)$ depends on conformal time
alone (spatial homogeneity), the optical index built from the conformal Killing
vector,
\begin{equation}
    n(\eta,r)=\frac1v=\frac{\Ehat}{\sqrt{\Ehat^2-a(\eta)^2}},
    \label{eq:flrwindex}
\end{equation}
is \emph{position-independent} at fixed $\eta$. The argument is then elementary and does
not even need Fermat's principle: the rail fixes the comoving speed to $v(\eta)$, so the
Euclidean arc length accumulated up to conformal time $\eta_f$ is
$\int_{\eta_0}^{\eta_f}v\,\dd\eta$, a strictly increasing function of $\eta_f$ alone.
Reaching the target therefore costs the least conformal time exactly when the comoving
path is the shortest, i.e.\ a straight line. The same curve minimizes $T_t$ and $T_\tau$
because on the rail \emph{every} clock is a function of $\eta$ alone: $\dd t=a\,\dd\eta$,
while the rail-fixed speed makes the Lorentz weight collapse,
\begin{equation}
   \frac{\dd\tau}{\dd\eta}=a\sqrt{1-v^2}=\frac{a^2}{\Ehat},
   \label{eq:flrw-dtau}
\end{equation}
so $t_f$ and $\tau_f$ are strictly increasing functions of $\eta_f$ and share its
minimizer. Hence
\begin{equation}
\begin{split}
    \arg\min T_t&=\arg\min T_\tau=\arg\min T_\eta\\
    &=\text{comoving straight line}.
\end{split}
    \label{eq:flrwdegen}
\end{equation}
The $t/\tau$ splitting that drives every later result is a gravitational-%
\emph{gradient} effect ($\partial_r f\neq0$); it is absent whenever the metric
functions depend on time only. FLRW is thus the degenerate base of the ladder:
it exhibits freezing at $a=\Ehat$ but, under perturbation, no branch structure
(figure~\ref{fig:flrw-kin}a,\,b).

\begin{figure}[tbp]
    \centering
    \begin{minipage}[t]{0.48\textwidth}\centering
        % ALT: Two stacked plots against cosmic time: four local-speed curves falling to
        %      zero, each ending at its own vertical dotted line; below, four decaying
        %      rail-redshift curves crossing a dashed unit line.
        \includegraphics[width=\textwidth]{Immagini/fig_flrw_cinematica.pdf}\\{\footnotesize (a)}
    \end{minipage}\hfill
    \begin{minipage}[t]{0.48\textwidth}\centering
        % ALT: Two travel-time curves against the amplitude of a sinusoidal path
        %      deformation, both smooth with their minimum at zero amplitude.
        \includegraphics[width=\textwidth]{Immagini/fig_flrw_variazionale.pdf}\\{\footnotesize (b)}
    \end{minipage}
    \caption{FLRW base case. (a)~Rail kinematics: local speed $v$ and rail redshift
    $E_{\rm phys}=\Ehat/a$, with freezing at $a=\Ehat$. (b)~Numerical perturbation check:
    sinusoidal deformations of the comoving straight line raise \emph{both} $T_t$ and
    $T_\tau$, consistent with the analytically proved branch degeneracy---a local check,
    not a global-minimality proof.}
    \label{fig:flrw-kin}\label{fig:flrw-var}
\end{figure}

\subsection{Arbitrary spatial curvature}
\label{sec:flrw-curvature}
Nothing above used spatial flatness. Write the slices with arbitrary curvature
$k=+1,0,-1$ in the conformally static form
\begin{equation}
   g=a(\eta)^2\Big[-\dd\eta^2+\dd\chi^2+S_k(\chi)^2\dd\Omega^2\Big],
   \qquad S_k(\chi)=\sin\chi,\ \chi,\ \sinh\chi ,
   \label{eq:flrw-k}
\end{equation}
and keep the same conformal selector $W=\partial_\eta$. Then $W\cdot W=g_{\eta\eta}=-a^2$,
so
\begin{equation}
   |W|=a(\eta)\qquad\text{for every }k :
   \label{eq:Wnorm-k}
\end{equation}
the conformal factor multiplies the whole metric and $W$ points along the conformal time,
so the spatial curvature never reaches the norm of the selector. By
Lemma~\ref{lem:indicatrix} the indicatrix is therefore the same centred sphere of speed
$v(\eta)=\sqrt{1-a^2/\Ehat^2}$ as in the flat case, the freezing surface is again
$a=\Ehat$, and the clock relation~\eqref{eq:flrw-dtau} is unchanged. Every clock remains a
strictly increasing function of $\eta$ alone, so the branch degeneracy~\eqref{eq:flrwdegen}
persists and the minimiser is the shortest comoving path. Hence for all three curvatures
\begin{equation}
   \text{brachistochrone}=\text{spatial geodesic of the slice},
   \qquad
   \int_{\eta_0}^{\eta_f}v(\eta)\,\dd\eta=d_\gamma(x_0,x_f),
   \label{eq:flrw-geodesic}
\end{equation}
with $d_\gamma$ the geodesic distance of the constant-curvature $3$-metric: curvature
enters \emph{only} through that distance. The comoving shapes differ---great-circle arc,
straight line, hyperbolic geodesic---as shown in
figure~\ref{fig:flrw-curvature}(a).

The closed case is not merely cosmetic: its slices are compact, and it supplies the
simplest concrete instance of the obstruction that Proposition~\ref{prop:hjb} can only
describe in the abstract. Two points of $S^3$ at geodesic distance $\pi$---antipodes---are
joined by a whole $S^2$ of minimising geodesics, a textbook Maxwell set and cut locus.
Whether the rail can reach it is quantitative: with $\Delta_{\max}$ the comoving distance
accumulated up to freezing, the antipode is attainable iff $\Delta_{\max}\ge\pi$ in units
of the comoving curvature radius $R_{\rm c}$. For de Sitter,
$\Delta_{\max}=\big[\sqrt{\Ehat^2-1}+\arcsin(1/\Ehat)-\pi/2\big]/(\Ehat H)$, so the
threshold is a condition on $HR_{\rm c}$ alone
(figure~\ref{fig:flrw-curvature}(b)): a closed universe whose curvature radius is large
compared with the Hubble radius keeps the minimiser unique, while a small enough one lets
the extremal reach its own cut locus, where global optimality fails even though every
competing arc is a perfectly good Pontryagin extremal. The open case $k=-1$ adds no such
feature: the slices are non-compact, geodesics between two points are unique, and the
structure is that of the flat case throughout.

\begin{figure}[tbp]
    \centering
    % ALT: Panel (a): three comoving curves leaving a common height above a marked point O,
    %      the flat one horizontal, the closed one bending toward O, the open one away.
    %      Panel (b): three saturating curves of reachable distance against rail energy,
    %      with a horizontal line at pi and the region above it shaded.
    \includegraphics[width=\textwidth,height=0.68\textheight,keepaspectratio]{Immagini/fig_flrw_curvatura.pdf}
    \caption{FLRW at arbitrary spatial curvature. (a)~Comoving brachistochrones with a
    common closest approach $\chi_{\min}$ to a reference point $O$, drawn in the geodesic
    polar map $(\chi,\varphi)$ of the slice: the flat geodesic is the horizontal line at
    height $\chi_{\min}$, the closed one bends toward $O$ and the open one away from it.
    Only the half $\varphi\ge0$ is shown; the curves are symmetric about the closest
    approach. (b)~de Sitter reach before freezing, $\Delta_{\max}$ in comoving curvature
    radii, against the antipodal distance $\pi$ of a closed slice. Above that line the cut
    locus is attainable and the minimiser stops being unique---the Maxwell obstruction of
    Proposition~\ref{prop:hjb} realised in the simplest member of the ladder.}
    \label{fig:flrw-curvature}
\end{figure}


% ======================================================================
\section{Vaidya: dynamical black hole and the Kodama rail}
\label{sec:vaidya}
% ----------------------------------------------------------------------
The ingoing Vaidya metric $\dd s^2=-f\,\dd v^2+2\,\dd v\,\dd r+r^2\dd\Omega^2$,
$f=1-2m(v)/r$ \cite{Vaidya1951}, has no timelike
Killing vector; for $m'(v)>0$ its trapping structure is that of a dynamical horizon
\cite{AshtekarKrishnan2003,Booth2005,FaraoniBook}. The restriction is not
decorative: Ashtekar and Krishnan make \emph{space-like} part of the definition of a
dynamical horizon \cite[Def.~1]{AshtekarKrishnan2003}, and the tube $r=2m(v)$ has
tangent $T=\partial_v+2m'(v)\partial_r$ with
\begin{equation}
    g(T,T)\big|_{r=2m(v)}=4\,m'(v) ,
    \label{eq:tube-signature}
\end{equation}
so it is spacelike, null or timelike exactly as $m'(v)$ is positive, zero or
negative. Accretion gives a dynamical horizon; the evaporating sign gives the
timelike tube they treat separately, which is the analytic continuation referred to
in section~\ref{sec:intro} and not a case we claim. It is worth noting which sign
the classical sources treat, since it is the opposite one. Vaidya's paper is
titled for it---``the gravitational field of a \emph{radiating} star'', a star
``supposed to start radiating at time $t_0$'' \cite[\S3]{Vaidya1951}---and
Lindquist, Schwartz and Misner analyse the same emitting case in the outgoing form
$\dd s^2=-[1-2m(u)/r]\dd u^2-2\,\dd u\,\dd r+r^2\dd\Omega^2$ in retarded time,
where positivity of the radiated energy density forces $\dd m/\dd u\le0$
\cite[\S III]{LindquistSchwartzMisner1965}. The ingoing form used here is the
advanced-time counterpart, describing accretion; it is standard, and standardly
called Vaidya's, but it is the time-reverse of what either paper computes, and the
mass-rate sign that distinguishes them is exactly the one the horizon signature
turns on.
Equation~\eqref{eq:tube-signature} is verified in
\texttt{../paper2/verification/verify\_vaidya\_tube.wls}. The Kodama vector
\cite{Kodama1980,Hayward1996,AbreuVisser2010}
\begin{equation}
    K^a=\varepsilon^{ab}\nabla_b r = \partial_v ,\qquad
    \nabla_a K^a=0,
    \label{eq:kodama}
\end{equation}
is divergence-free without any Killing symmetry, and its charge is the
Misner--Sharp mass $m(v)$ \cite{MisnerSharp1964,HaywardMukohyama1999}---in
Hayward's words, ``the conserved Kodama current has charge $E$'', with
$E=\tfrac12r\big[1-g^{-1}(\dd r,\dd r)\big]$ \cite[Eq.~(4)]{Hayward1996}, which for
ingoing Vaidya evaluates to $m(v)$. His trapping criterion $E\gtrless r/2$ then
places the marginal sphere at $r=2m(v)$, as used throughout.

Two remarks on the normalisation and the domain, both from the original.
Kodama's own definition \cite[Eq.~(3.6)]{Kodama1980} carries a different
normalisation from the modern $K^a=\varepsilon^{ab}\nabla_br$ used here, and he
notes that his $K$ ``coincides with the Killing vector in its direction but not
in its norm''. With the present normalisation the agreement is exact: for the
ingoing chart one finds $K=\partial_v$ and $g(K,K)=-f$ identically, so in the
static limit $K$ reduces to $\xi=\partial_t$ in norm as well as direction. And
Kodama states there that $K$ is timelike \emph{outside the apparent horizon}
$\Phi=0$. We follow the numerical-relativity usage of that term throughout---the
outermost marginally trapped surface on a slice, here $r=2m(v)$---rather than the
Hawking--Ellis definition as the boundary of the trapped region; the two are
routinely conflated, and both Ashtekar and Krishnan
\cite[footnote~1]{AshtekarKrishnan2003} and Booth
\cite[\S1]{Booth2005} take care to separate them. That is precisely the boundary at which the compact-control problem
ends (Paper~II, Lemma on compactness of the rail slice): the selector ceases to
be timelike exactly where the original construction says it does, which is a
third independent route to the same surface.

\subsection{Kodama energy is the rail invariant}
\label{sec:kodama-energy}
With the canonical selector $K=\partial_v$, the rail constraint~\eqref{eq:rail} is
exactly $-u\cdot K=-u_v=\Ehat$, hence
\begin{equation}
    -\,u\cdot K \;=\; \Ehat \quad\text{identically along the brachistochrone.}
    \label{eq:kodama-conserved}
\end{equation}
This is \emph{controlled} conservation, not Noether: for a geodesic
$-u\cdot K$ would drift at rate $-m'(u^v)^2/r\propto m'$, exactly the
right-hand side of equation~\eqref{eq:cost} for $W=K$. Only the unnormalized
$K=\partial_v$ has zero control cost in the static limit; the normalized
$\hat K=K/\sqrt f$ does not (discriminating test), which is why the Kodama
vector, not a unit selector, is the canonical choice.

\paragraph{Indicatrix and Hamiltonians.} With $w\equiv\Ehat^2-f$ the constraint
$-u_v=\Ehat$ turns the admissible velocities at fixed $(v,r)$ into an ellipse,
\begin{equation}
    r'=(f-\Ehat^2)+\Ehat\sqrt w\,\cos\theta,\qquad
    \varphi'=\frac{\sqrt w}{r}\sin\theta,
    \label{eq:vaidya-ind}
\end{equation}
a Zermelo problem with ingoing radial wind $f-\Ehat^2<0$. Maximizing
$p\cdot\dot x$ over $\theta$ at fixed $p_\varphi=J$ gives the closed-form branch
Hamiltonians
\begin{align}
    H_v&=p_r(f-\Ehat^2)-1+\sqrt w\,\sqrt{\Ehat^2 p_r^2+J^2/r^2},
    \label{eq:Hv}\\
    H_\tau&=p_r(f-\Ehat^2)-\Ehat
        +\sqrt w\,\sqrt{(\Ehat p_r+1)^2+J^2/r^2}.
    \label{eq:Htau-vaidya}
\end{align}
Both branches are built from the \emph{same} support function of the same
indicatrix~\eqref{eq:vaidya-ind},
\begin{equation}
   h(p_r,J)=p_r\,(f-\Ehat^2)+\sqrt w\,\sqrt{\Ehat^2p_r^2+J^2/r^2},
   \label{eq:vaidya-support}
\end{equation}
and differ only in the cost subtracted from it: the $v$-branch has unit cost,
$H_v=h(p_r,J)-1$, whereas the proper-time cost $\dd\tau/\dd v=(f-r')/\Ehat$ is itself
control-dependent and is absorbed by shifting the radial momentum,
\begin{equation}
   H_\tau(p_r,J)=h\!\left(p_r+\tfrac1{\Ehat},\,J\right)-\frac{f}{\Ehat} .
   \label{eq:tau-from-v}
\end{equation}
Carrying out the shift reproduces~\eqref{eq:Htau-vaidya}: the constant $-\Ehat$ there is
the \emph{result} of combining $-f/\Ehat$ with the $(f-\Ehat^2)/\Ehat$ generated by the
shift, not the cost itself. Non-stationarity is the explicit dependence
$\dd H/\dd v=\partial_v H\propto m'(v)$; the \emph{endpoint-sensitive} (arising from the
two-point boundary-value problem with a free future arrival---a costate memory, not a
causal effect) momentum $p_v=-H$
satisfies $p_v(r_1)=0$ at the free arrival radius yet $p_v(r)\neq0$ along the
path, so the turning point \emph{anticipates future accretion}
(figure~\ref{fig:vaidya-traj}).

\paragraph{Branch costate vs mechanical angular momentum.} The conserved $J=p_\varphi$
is the Noether \emph{costate} of the reduced control Hamiltonian---the canonical
momentum conjugate to $\varphi$ in the arrival-time problem---\emph{not} the spacetime
angular momentum $L_{\rm mech}=g(u,\partial_\varphi)$. Although $\partial_\varphi$ is a
Killing vector, Noether's theorem for the \emph{control} symmetry returns the control
costate, not the mechanical momentum: along a forced (rail-constrained, accelerated)
worldline $L_{\rm mech}$ is \emph{not} conserved, since
$\dd L_{\rm mech}/\dd\tau=a\cdot\partial_\varphi\neq0$. The two are therefore distinct
objects. On the native Vaidya $v$-branch shell the Legendre map gives the exact
relation
\begin{equation}
    L_{\rm mech}=r^2 u^\varphi=\frac{\Ehat^2-f(r)}{\Ehat}\,J,
    \label{eq:costate-mech}
\end{equation}
which varies with $r$ while $J$ stays constant (they agree only at the isolated radius
$f=\Ehat^2-\Ehat$). This follows from the indicatrix in two lines: at the maximizer
$(\cos\theta^\star,\sin\theta^\star)\propto(\Ehat p_r,\,J/r)$, so with
$N=\sqrt{\Ehat^2p_r^2+J^2/r^2}$ one gets
$L_{\rm mech}=\sqrt w\,J/[\Ehat(N-\sqrt w\,p_r)]$, while the shell $H_v=0$ reads
$\sqrt w\,N=1+p_rw$, i.e.\ $N-\sqrt w\,p_r=1/\sqrt w$; substituting
gives~\eqref{eq:costate-mech}. The Legendre
map is branch-specific, but on all Vaidya branches the conserved axial control costate is
$J$ and the slow deformation is generated solely by $m\partial_m$; we reserve ``angular
momentum'' for the mechanical $g(u,\partial_\varphi)$ and call $J$ the branch costate,
$L_{\rm mech}$ being a derived, non-conserved quantity.

\begin{figure}[ht]
    \centering
    % ALT: Two stacked plots against radius: above, on a logarithmic axis, a nearly flat
    %      residual trace below 1e-15; below, a smooth nowhere-vanishing drift curve.
    \includegraphics[width=\textwidth,height=0.68\textheight,keepaspectratio]{Immagini/fig_kodama_conservazione.pdf}
    \caption{Kodama energy along the $\tau$-brachistochrone. Top: the
    \emph{residual} $|{-}u\cdot K-\Ehat|$ on a \emph{log} scale, held below
    $6.7\times10^{-16}$---an \emph{actively controlled} (rail-maintained) invariant,
    not a Noether constant. Bottom: the free-particle Kodama drift $\propto m'$ that
    the rail cancels (the missing symmetry).}
    \label{fig:kodama}
\end{figure}

\paragraph{On-shell drift versus off-shell shell departure.} For the Vaidya family the frozen orbit shape
$\mathrm{d}\varphi/\mathrm{d}r=F(r;\Ehat,J)$ depends on the mass modulus $m$, and its
on-shell drift is generated by $m\partial_m$. Conservation of the Kodama rail charge
$\Ehat$ does \emph{not} constrain the non-autonomous trajectory to the instantaneous
$H=0$ shell---the costate equation is independent information---so the \emph{complete}
first-order response adds the $p_v$-driven off-shell contribution
[equation~\eqref{eq:adiab-exact}, section~\ref{sec:vaidya-adiabatic},
Figure~\ref{fig:vaidya-offshell}].


\subsection{Phenomenology: penetration, timing, bounce}
\label{sec:vaidya-pheno}
Three features distinguish the dynamical hole from its frozen-mass
Schwarzschild limit.

\emph{Penetration.} The $\tau$-brachistochrone reaches the apparent horizon
$r=2m(v)$ only if $|J|$ lies below a threshold $\Jpen(v_0)$; accretion ($m'>0$)
opens a \emph{finite} window $\Jpen>0$, whereas the negative-rate continuation
($m'<0$ in the \emph{ingoing} metric---a formal continuation; physical evaporation is
the \emph{outgoing} metric, section~\ref{sec:vaidya-plunge}) closes it to zero
measure---an accretion/(negative-rate) asymmetry with no static analogue. The
finite capture interval $J\in(0,\Jpen]$ is the analogue of the Kerr $t$-branch
dichotomy, generated here by the \emph{dynamics}: grazing orbits linger near
periapsis while the growing horizon closes the gap.

\emph{Bounce.} The $r$-parametrization degenerates at periapsis
($\dd v/\dd r\to\infty$); reparametrizing by $v$ makes the Euler--Lagrange
problem regular at the turning point $r'=0$, so the trajectory is integrated
straight \emph{through} the bounce---the same extremal, via the exact change of
variable $(f v'-1)\dd r=(f-r')\dd v$. The $v$-flow reaches the periapsis exactly
where the $r$-flow diverges ($r_{\min}=3.0105$ at $v=7.29$ for $\mu\equiv m'=0.01$),
and the endpoint-sensitive memory $p_v(r)\propto m'$ measures the departure from the
frozen-mass Schwarzschild curve (figure~\ref{fig:vaidya-bounce}).

\emph{Timing and the orbit--horizon race.} Because $m$ grows along the path,
$r_{\min}$ depends on \emph{when} the particle arrives, the epoch $v_1$---a
genuinely non-stationary effect absent in Kerr. The absolute and
relative penetration have \emph{opposite} winners
(figure~\ref{fig:vaidya-timing}): under the negative-rate branch ($m'<0$) the absolute
periapsis descends but the horizon $2m(v)$ retreats faster, so the relative penetration
$r_{\min}/2m(v_{\rm peri})$ \emph{worsens} for later arrivals; under accretion the
periapsis rises yet the horizon grows faster, so late arrivals graze the relative
horizon. The orbit--horizon chase has reversed outcomes measured in $r$ versus in
horizon units.

\begin{figure}[ht]
    \centering
    % ALT: Two nearly coincident equatorial orbit shapes separating inward, and a second
    %      panel showing the costate p_v against radius, zero at the outer endpoint and
    %      nonzero along the inward path.
    \includegraphics[width=\textwidth,height=0.68\textheight,keepaspectratio]{Immagini/fig_vaidya_traiettorie.pdf}
    \caption{$\tau$-brachistochrones for Schwarzschild and Vaidya, and $p_v(r)$,
    at $E=1.2$, $J=1.3$, $r_1=10$, arrival time $v_1=40$, with
    $m(v)=1+\mu v$ and $\mu\in\{0,0.01\}$. Transversality gives
    $p_v(r_1)=0$, while Vaidya develops non-autonomous memory $p_v\propto m'$.}
    \label{fig:vaidya-traj}
\end{figure}


\subsection{Plunge-radius law and the absence of inversion in the frozen-clock comparison}
\label{sec:vaidya-plunge}
The turning-point law has the static Schwarzschild form evaluated at the mass
at periapsis, $f_*=1-2m(v_{\rm peri})/r_{\min}$:
\begin{align}
    \tau:&\quad (\Ehat^2-f_*)J^2 = f_* r_{\min}^2 ,\\
    t:&\quad (\Ehat^2-f_*)J^2 = \Ehat^2 r_{\min}^2/f_* ,
    \label{eq:vaidya-turn}
\end{align}
plus a genuine $O(m')$ correction from the endpoint-sensitive costate memory $p_v$. Both
lines follow from Lemma~\ref{lem:indicatrix} in the frozen limit. Parametrizing by the
Killing time, the rail $-u\cdot\partial_t=\Ehat$ gives $\dd t/\dd\tau=\Ehat/f$, and
substituting into $g(u,u)=-1$ turns the admissible velocities
$(V^r,V^\varphi)=(\dd r/\dd t,\ r\,\dd\varphi/\dd t)$ into the \emph{centred} ellipse
\begin{equation}
   \frac{(V^r)^2}{f^2}+\frac{(V^\varphi)^2}{f}=\frac{w}{\Ehat^2},
   \qquad
   a_r=\frac{f\sqrt w}{\Ehat},\quad a_\varphi=\frac{\sqrt{f\,}\sqrt w}{\Ehat},
   \label{eq:frozen-ind-t}
\end{equation}
centred because the frozen metric is static: there is no wind, only eccentricity. Its
support function is $h=\sqrt{a_r^2p_r^2+a_\varphi^2J^2/r^2}$, and the two branches differ
only in the cost subtracted from it---unity for $t$, and $\dd\tau/\dd t=f/\Ehat$ for
$\tau$, which in this parametrization is control-\emph{independent}, unlike the $v$-branch
cost of~\eqref{eq:tau-from-v}. Setting $p_r=0$ at the turning point in
$h=1$ and in $h=f/\Ehat$ returns the two lines of~\eqref{eq:vaidya-turn} respectively. The
static ratio $J^2_t/J^2_\tau=\Ehat^2/f^2>1$ gives $r_{\min}^t<r_{\min}^\tau$
(``$t$ deeper''), a \emph{branch-intrinsic} ordering set by $\Ehat/f>1$ alone---it
does not depend on the sign of $m'$ nor on the ingoing/outgoing character of the
flux. \emph{Physical} evaporation is the \emph{outgoing} Vaidya metric
$\dd s^2=-f\dd u^2-2\dd u\,\dd r$ (retarded time $u$), not the ingoing metric with
$m'<0$. The depth ordering $r_{\min}^t<r_{\min}^\tau$ is \emph{branch-intrinsic}: it
follows from $\Ehat/f>1$ alone and is therefore independent of the sign of $m'$ and of
the ingoing/outgoing character of the flux, so to this (static, frozen-clock) order it
is unchanged under evaporation. The apparent linear-model ``inversion'' from flipping
$m'<0$ in the \emph{ingoing} model is an unphysical extrapolation to $m\le0$. We do
\emph{not} here solve the full outgoing optimal-control problem with its own
Hamiltonians, costates and transversality: the advanced/retarded frozen-clock asymmetry
[equation~\eqref{eq:vaidya-asymmetry}] is a clock decomposition ($v\!\leftrightarrow\!u$,
the tortoise $r_*$ changing sign) of the frozen first-order correction, not a complete
outgoing boundary-value solution. Accordingly we state the absence of an evaporative
inversion as a property of the branch-intrinsic static ordering together with this
frozen clock decomposition, and leave the full outgoing BVP---and a genuine outgoing
trajectory to accompany figure~\ref{fig:vaidya-noinv}---to future work.

\begin{figure}[ht]
    \centering
    % ALT: Two decreasing plunge-radius curves against the mass-rate parameter with a third
    %      curve for their difference; the coordinate-time curve lies below throughout and
    %      the gap widens toward positive rate.
    \includegraphics[width=\textwidth,height=0.68\textheight,keepaspectratio]{Immagini/fig_vaidya_plunge_t_tau.pdf}
    \caption{Plunge radii $r_{\min}^{t,\tau}$ and their gap versus the ingoing
    mass-rate parameter $\mu=m'$ for $E=1.2$, $J=8$, $r_1=15$, $v_1=40$ and
    $m(v)=1+\mu v$. The $t$ branch remains deeper and accretion widens the gap.}
    \label{fig:vaidya-plunge}
\end{figure}

\begin{figure}[ht]
    \centering
    % ALT: The branch gap against the mass-rate parameter, crossing zero only inside a
    %      shaded unphysical negative-mass band, with a second curve approaching zero from
    %      above without crossing.
    \includegraphics[width=\textwidth,height=0.68\textheight,keepaspectratio]{Immagini/fig_vaidya_no_inversione_evaporazione.pdf}
    \caption{Absence of inversion in the ingoing, positive-mass model: the
    linear-model sign change lies entirely in the unphysical $m<0$ region, and with
    $m=e^{-\lambda v}>0$ the branch gap shrinks but does not invert. This is a
    frozen-clock branch comparison at $E=1.2$, $J=8$, $r_1=15$, $v_1=40$, not a
    solution of the outgoing-Vaidya boundary-value problem (which remains open).}
    \label{fig:vaidya-noinv}
\end{figure}

\begin{figure}[htbp]
    \centering
    % ALT: Radius against advanced time through the turning point, a smooth rounded minimum
    %      continuing outward, overlaid on the radius-parametrised flow which terminates at
    %      that minimum.
    \includegraphics[width=\textwidth,height=0.68\textheight,keepaspectratio]{Immagini/fig_vaidya_bounce.pdf}
    \caption{Bounce through periapsis in the $v$-parametrization: the
    Euler--Lagrange flow is regular at $r'=0$ and integrates through the turning
    point, reaching $r_{\min}$ exactly where the $r$-parametrization diverges.
    Parameters are $E=1.2$, $J=1.3$, $r_1=10$, $v_1=40$, and
    $m(v)=1+\mu v$ with $\mu\in\{-0.01,0,0.01\}$.}
    \label{fig:vaidya-bounce}
\end{figure}

\begin{figure}[htbp]
    \centering
    % ALT: Two plots against arrival epoch, each with three curves for the three mass rates:
    %      periapsis radius, and the same radius divided by the horizon radius; the ordering
    %      of the curves reverses between the two.
    \includegraphics[width=\textwidth,height=0.68\textheight,keepaspectratio]{Immagini/fig_vaidya_timing.pdf}
    \caption{Timing window: $r_{\min}$ depends on the arrival epoch $v_1$. The
    orbit--horizon race has opposite winners in absolute ($r_{\min}$) and relative
    ($r_{\min}/2m$) terms for accretion versus the negative-rate branch. Here
    $E=1.2$, $J=1.3$, $r_1=10$, $v_1\in[25,70]$, and
    $\mu=m'\in\{-0.01,0,0.01\}$ in the ingoing model. For the earliest arrivals the
    periapsis falls at $v\simeq-7$, so the linear mass model is evaluated slightly before
    $v=0$; it stays positive there ($m\ge0.93$), and no $m\le0$ extrapolation occurs.}
    \label{fig:vaidya-timing}
\end{figure}

\subsection{Closed-form adiabatic correction and the horizon dilogarithm}
\label{sec:vaidya-adiabatic}
When $m(v)$ varies slowly we apply the general frozen-coefficient reduction of
Lemma~\ref{lem:reduction} and the first-order response of Theorem~\ref{thm:firstorder} to the Vaidya mass modulus, with
$\dot m=\dd m/\dd v$ the small rate and $\varepsilon=M\dot m/m$ its dimensionless
bookkeeping parameter. In the spectral formulae from this point onward we abbreviate
$E\equiv\Ehat$ and define the single polynomial factor
$D_E\equiv(E^2-1)r+2m$. At a fixed reference slice the two normalizations satisfy
$\delta\varphi/\varepsilon=(m/M)\,\delta\varphi/\dot m$. The frozen family
is the \emph{Schwarzschild} brachistochrone---the $a\to0$ limit of the
Thakurta--Kerr forms---on the sextic
$S=r(r-2m)D_E\,[r^2(r-2m)-J^2 D_E]$. Concretely it is the \emph{proper-time} branch
$H_\tau$ of~\eqref{eq:Htau-vaidya}, on the orientation of increasing $r$ that renders
$\sqrt S$ positive: solving $H_\tau=0$ for $p_r$ and forming
$\varphi'=H_{\tau,J}/H_{\tau,p_r}$, $\dd v/\dd r=1/H_{\tau,p_r}$ and $\dd\tau/\dd r$
returns exactly the three spectral expressions used below,
$\varphi'=J D_E/\sqrt S$, equation~\eqref{eq:vaidya-clock} and
$\dd\tau/\dd r=r^2(r-2m)/\sqrt S$ (checked to $10^{-25}$ at several radii). Notably this
sextic keeps six distinct roots: the constrained brachistochrone is genus \emph{two} even in
Schwarzschild (only the geodesics are elliptic; the genus drop occurs at the algebraic
degeneration value $|J|=J_{\rm deg}$). The mass derivative reduces exactly, the horizon factor cancelling,
\begin{equation}
    \partial_m F=\frac{N_m}{S^{3/2}},\quad N_m=S\,\partial_m K-\tfrac12 K\,\partial_m S,
    \qquad
    \partial_m\varphi_0=\frac{\mathcal{A}^m}{\sqrt S}+\sum_{k=0}^4 c_k^m U_k,
    \label{eq:dm-reduction}
\end{equation}
with $U_k=\int r^k/\sqrt S\,\dd r$ and \emph{algebraic} constants $c_k^m$ (verified
to $10^{-15}$). The clock is now the \emph{advanced time}: from $v=t+r_*$,
\begin{equation}
    \frac{\dd v}{\dd r}=\frac{E}{f}\frac{\dd\tau}{\dd r}+\frac1f
    =\frac{E\,r^3}{\sqrt S}+\frac{r}{r-2m}
    \;\Rightarrow\;
    v=E\,U_3+r+2m\ln(r-2m),
    \label{eq:vaidya-clock}
\end{equation}
the tortoise term injecting a \emph{horizon logarithm}. This is the advanced-time
(native) branch; the \emph{proper-time} branch instead uses the second-kind clock
$\dd\tau/\dd r=r^2(r-2m)/\sqrt S$, i.e.\ $T_\tau=U_3-2m\,U_2$, $b=(0,0,-2m,1,0)$ (no
horizon log), and closes by the very same
universal on-shell assembly established here (specialized to the Vaidya clock)---$\delta\varphi_\tau/\dot m=\tfrac12\sum_{k<j}Q_{kj}
W_{kj}-T_\tau(\mathcal{A}^m/\sqrt S+\tfrac12\sum c_k^m U_k)+\sum_j b_j\mathcal{I}_j$ with
$Q_{kj}=c_k^m b_j-c_j^m b_k$, all coefficients algebraic (verified to
$10^{-14}$). The two canonical clocks of the ingoing problem are therefore the
\emph{advanced null time} $v$ (adapted to the Kodama selector $K=\partial_v$, the
native evolution parameter) and the \emph{proper time} $\tau$; the coordinate time
$t$ that enters the static depth comparison~\eqref{eq:vaidya-turn} is the
\emph{frozen} Schwarzschild-like time, a comparison device valid at frozen order and
\emph{not} a canonical non-autonomous clock---once $m=m(v)$ the tortoise
$r_*=r+2m\ln(r-2m)$ acquires $\dd m$ terms and no Killing time survives. The retarded
time $u=t-r_*$ of the outgoing metric is a \emph{separate} outgoing-Vaidya problem
(its own Hamiltonians, costates, and transversality), not a third clock of the ingoing
family. The advanced-time first-order correction
$\delta\varphi=\dot m\!\int\!\partial_m F\,v\,\dd r$ then integrates by parts into a
fully explicit closed form (verified to $9\times10^{-14}$, coefficients given),
\begin{equation}
    \frac{\delta\varphi}{\dot m}
    =A_m\,v-E\!\int\! A_m\,\frac{r^3}{\sqrt S}\,\dd r
    -\int\! A_m\,\frac{r}{r-2m}\,\dd r,\qquad A_m\equiv\partial_m\varphi_0,
    \label{eq:vaidya-deltaphi}
\end{equation}
This is the \emph{on-shell} component of the Vaidya first-order correction (the
mass-drift of the frozen family).

\Ithm{5}
\begin{theorem}[complete first-order response]\label{thm:firstorder}
Assume (H) and let the background depend on a single slow parameter with dimensionless
rate $\varepsilon$. Linearizing the extended Hamiltonian $\mathcal H=p_\chi+H=0$ about the
frozen extremal, in the original canonical variables $(r,p_r)$ with $\lambda$ the elapsed
clock measured from the free-arrival endpoint, the momentum-shell shift is
$\delta p_r=(\lambda\,\Theta H-S_{\mathcal D})/H_{p_r}$ and the shape correction is
\begin{equation}
    \begin{aligned}
    &\frac{\delta\varphi}{\varepsilon}=\int\!\Big[\frac{\partial_{p_r}G}{H_{p_r}}\,
    \big(\lambda\,\Theta H-S_{\mathcal D}\big)-\lambda\,\Theta G\Big]\dd r,\qquad
    G=\frac{H_J}{H_{p_r}}=\varphi',\\[2pt]
    &S_{\mathcal D}(\lambda)=\!\int_0^{\lambda}\!\mathcal D H\,\dd\lambda',
    \end{aligned}
    \label{eq:adiab-exact}
\end{equation}
which closes the correction to $O(\varepsilon^2)$, uniformly on any compact sub-arc on
which $|H_{p_r}|$ is bounded away from zero. Here $\Theta$ is the Euler operator
generating the frozen-family drift and $\mathcal D$ is the drift generator, equal to
$\Theta$ unless the slow parameter rescales the canonical momenta. For single-parameter
Vaidya, $\mathcal D=\Theta=m\partial_m$ and the terminally-anchored source collapses to
the universal boundary form
\begin{equation}
    S_{\mathcal D}=[\,r\,p_r\,]-\lambda .
    \label{eq:universal-boundary}
\end{equation}
\end{theorem}

\begin{proof}
The derivation is metric-independent and is carried out, then closed explicitly for the
two Vaidya branches, in~\ref{app:firstorder}; the boundary form~\eqref{eq:universal-boundary}
follows there from the Randers--Euler identity together with the degree-zero
self-similarity of the frozen Hamiltonian, and the $O(\varepsilon^2)$ remainder from a
Gr\"onwall estimate on the sub-arc.
\end{proof}

As in Thakurta--Kerr, the interior costate does not
vanish---$p_v=-H_v=O(\dot m)$ with $p_v(r_1)=0$ only at the free arrival---so the
\emph{complete} first-order term adds the off-shell piece given by the Vaidya analogue
of equation~\eqref{eq:adiab-exact}, with the single-parameter operator $\Theta=m\partial_m$
and $S_{\mathcal D}=\int_0^\lambda\Theta H_v\,\dd\lambda'$ terminally anchored; the extra term is
$-\!\int(\partial_{p_r}G/H_{p_r})\,S_{\mathcal D}\,\dd r$, a length-two iterated integral of the
same horizon-dilogarithm class. Applying the identical canonical perturbation theory
to $H_v$ closes the Vaidya first order to $O(\dot m^2)$ (residual slope $\approx2$;
\texttt{VaidyaMetric/vaidya\_first\_order\_offshell.py}). We display the on-shell
component in closed form below;
its two integrals expand into named building blocks with algebraic coefficients.
Writing $A_m=\mathcal{A}^m/\sqrt S-\alpha_0+\sum_k c_k^m U_k$ with the constant
$\alpha_0=\mathcal{A}^m(r_0)/\sqrt{S(r_0)}$, integration by parts in each factor gives
the \emph{fully assembled} closed form (verified to $2\times10^{-13}$)
\begin{equation}
\begin{aligned}
    \frac{\delta\varphi}{\dot m}
    ={}&A_m\,v
    -E\Big[\mathcal{I}_{\rm poly}-\alpha_0 U_3
        +\tfrac12\sum_{k=0}^4 c_k^m\big(U_kU_3-W_{3k}\big)\Big]\\
    &{}-\Big[\mathcal{L}_{2m}-\alpha_0\,\big(v-E\,U_3\big)\Big|_{r_0}^{\,r}
        +\sum_{k=0}^4 c_k^m\big(r\,U_k-U_{k+1}+2m\,\mathcal{D}_k\big)\Big],
\end{aligned}
    \label{eq:vaidya-full}
\end{equation}
All blocks below are normalized to \emph{vanish at the lower limit} $r_0$, i.e.\ every
integral is $\int_{r_0}^{r}$; this is what makes the boundary term of the integration by
parts reduce to $A_m(r)\,v(r)$, since $A_m(r_0)=0$ by the choice of $\alpha_0$. The one
elementary primitive that is not of this form is the tortoise term, which is why it is
displayed with an explicit bracket $\big|_{r_0}^{\,r}$: evaluating it at the upper limit
alone would spoil the identity by the constant
$\alpha_0\,[r_0+2m\ln(r_0-2m)]$. With that convention the assembled
right-hand side reproduces the quadrature of~\eqref{eq:vaidya-deltaphi} to $10^{-39}$ at
$m=1$, $E=7/5$, $J=5/2$, $r_0=5$, $r=9$ (forty-digit arithmetic).

The named building blocks are: the abelian integrals $U_k=\int r^k/\sqrt S\,\dd r$;
the elementary rational-log primitive
$\mathcal{I}_{\rm poly}=\int\mathcal{A}^m r^3/S\,\dd r$; the genus-two
polylogarithms $W_{3k}=\int(U_3\dd U_k-U_k\dd U_3)$ \emph{at infinity};
the \emph{second}-kind horizon integral
$\mathcal{L}_{2m}=\int\mathcal{A}^m r/[(r-2m)\sqrt S]\,\dd r$---since
$\mathcal{A}^m(2m)\ne0$ its integrand behaves as $(r-2m)^{-3/2}$, a double pole with
\emph{no} residue at the horizon Weierstrass point, closed by the Kleinian
$\zeta$/quasi-periods~\cite{BuchstaberEnolskii1997,BakerAbelian1897}, not by Fay; the
genuinely \emph{third}-kind object is the horizon letter $\dd r/(r-2m)$ (the tortoise
log, a simple pole of residue two at that Weierstrass point), which closes as a Fay
$\theta$-ratio~\cite{Fay1973} centred at $r=2m$; and
the \emph{horizon dilogarithms} $\mathcal{D}_k=\int U_k/(r-2m)\,\dd r$. Since $r=2m$
is a branch point of $S$, in the local uniformiser $t=\sqrt{r-2m}$ one has
$\dd r/(r-2m)=2\,\dd t/t$: a \emph{simple} pole of residue two, i.e.\ a genuine
third-kind letter $\dd\log$ anchored \emph{at the horizon}. On the compact
genus-two curve this letter carries the divisor $2\,(r{=}2m)-\infty_+-\infty_-$ (residues
summing to zero) and closes as a Fay prime-form $\theta$-ratio; one integration by parts
then closes $\mathcal{D}_k$ onto a \emph{single} weight-two letter (verified to
$10^{-12}$). Displayed horizon logarithms carry a fixed reference length,
$\ln(r-2m)\equiv\ln[(r-2m)/L]$, which cancels in every assembled result.
\begin{equation}
    \mathcal{D}_k=U_k\,\ln(r-2m)-\mathcal{G}_k,\qquad
    \mathcal{G}_k=\int\ln(r-2m)\,\frac{r^k}{\sqrt S}\,\dd r,
    \label{eq:horizon-dilog}
\end{equation}
where $\mathcal{G}_k$ is the horizon-anchored genus-two dilogarithm---the exact
counterpart of the Kerr weight-two piece $W_{jk}$, in the same
hyperelliptic-dilogarithm class (conjecturally irreducible to weight one, the
theta-divisor obstruction $\Theta=W_{g-1}$ preventing the genus-one factorization),
but with its $\dd\log$ letter sourced by the
advanced-time divergence at the horizon $r=2m$ rather than by a third-kind dipole at
infinity. At the genus-degeneration $|J|=J_{\rm deg}$ the sextic acquires a double root at $r=r_d$
and the genus drops to one: the $U_k$ reduce to elliptic integrals and $\mathcal{G}_k$
collapses to the tabulated elliptic (Bloch--Wigner/Zagier) dilogarithm
\cite{Bloch2000,Zagier1990} on $E:\,w^2=Q_4(r)$. For spherical Vaidya this locus is
\emph{algebraic} ($r_d<0$, a double root off the physical arc, not an external
separatrix, and not to be conflated with the dynamical penetration threshold
$\Jpen(v_0)$); the explicit Weierstrass $\sigma,\zeta$ reduction, with representative
numerical values, is given in \ref{app:degeneration}. Every coefficient is algebraic
(rational in $m,E,J$); the transcendence lives entirely in $W_{jk},\mathcal{G}_k$, so the
correction is reusable by direct evaluation once $S$ and the $c_k^m$ are fixed.

\paragraph{Explicit degeneration-family correction.} On the algebraic genus-degeneration
locus $J=J_{\rm deg}$ (a double root at $r_d<0$, not a physical separatrix) the correction
reduces to a finite basis of standard letters. The antiderivative of the source is a sum
of Weierstrass letters,
\begin{equation}
    G(z)=\int\partial_m F=C_0\,z+\sum_a\Big[b_1^a\ln\sigma(z-a)-b_2^a\zeta(z-a)
        -\tfrac12 b_3^a\wp(z-a)\Big],
    \label{eq:vaidya-sep-G}
\end{equation}
with every residue $b_n^a$ closed in algebraic form over the poles
$a\in\{\pm z_d,0,i\omega_{\rm im}\}$, while the weight-two blocks $W_{kj}$ are length-two
iterated (Brown--Levin elliptic-polylogarithm~\cite{BrownLevin2011}) integrals evaluated by
a convergent $q$-series. \ref{app:degeneration} gives the algebraic reduction,
Mittag--Leffler block structure and residue formulae. Construction of the period data and
numerical Kronecker--Eisenstein/$q$-series evaluation use the archived routine. Every
coefficient is an algebraic function of $(m,E)$, so
$\delta\varphi_\tau|_{\rm deg}$ is reconstructed on any curve by evaluating these formulae
(checked to $10^{-9}$ against the direct integral across $E\in[1.3,1.6]$). The minimal
dilogarithm-basis dimension and the irreducibility of the weight-two letters remain
\emph{conjectural}; a secondary numerical rank experiment is in the reproducibility
package.

\paragraph{Advanced/retarded frozen-clock decomposition.} At first order $\delta\varphi$ is
linear in $\dot m$, so na\"ively $\dot m\to-\dot m$ merely flips its sign. The
formal advanced/retarded frozen-clock asymmetry is subtler: accretion ($\dot m>0$) is
described by the \emph{ingoing} Vaidya model and its \emph{advanced} clock $v=t+r_*$,
whereas the \emph{retarded} clock $u=t-r_*$ ($\dd s^2=-f\dd u^2-2\dd u\,\dd r$) belongs to
the outgoing geometry. The frozen reduction $\partial_m F$ is identical
(it depends only on $S$); only the clock changes, and $v\leftrightarrow u$ flips the
sign of the tortoise $r_*$. Splitting the clock into its bulk and tortoise parts,
$v=E\,U_3+r_*$ with $r_*=(r-r_0)+2m\ln\!\frac{r-2m}{r_0-2m}$,
\begin{equation}
    \frac{\delta\varphi}{\dot m}\Big|_{\rm adv}=\mathcal{A}_\infty+\mathcal{B}_{\rm hor},
    \quad
    \frac{\delta\varphi}{\dot m}\Big|_{\rm ret}=\mathcal{A}_\infty-\mathcal{B}_{\rm hor},
    \quad
    \mathcal{A}_\infty=\!\int\!\partial_m F\,E\,U_3\,\dd r,\;\;
    \mathcal{B}_{\rm hor}=\!\int\!\partial_m F\,r_*\,\dd r,
    \label{eq:vaidya-asymmetry}
\end{equation}
so that the retarded frozen-clock branch is
$\delta\varphi_{\rm ret}=\dot m(\mathcal{A}_\infty-\mathcal{B}_{\rm hor})$; at equal
rate magnitude the \emph{formal advanced/retarded frozen-clock difference} is
$\delta\varphi_{\rm adv}-\delta\varphi_{\rm ret}=2|\dot m|\,\mathcal{A}_\infty$
(carried by the infinity bulk), while the \emph{deviation from exact antisymmetry} is
$\delta\varphi_{\rm adv}+\delta\varphi_{\rm ret}=2|\dot m|\,\mathcal{B}_{\rm hor}$,
carried \emph{entirely by the tortoise/horizon}: the horizon term $\mathcal{B}_{\rm hor}$
is exactly what breaks the na\"ive $\dot m\to-\dot m$ antisymmetry. This
advanced/retarded split is a \emph{frozen-clock} comparison at first order, isolating the
sign change of the tortoise contribution; it is \emph{not} a solution of the full
outgoing-Vaidya optimal-control boundary-value problem, whose Hamiltonians, costates and
transversality remain to be derived. Within
$\mathcal{B}_{\rm hor}$ the transcendental, direction-sensitive content is precisely
the horizon term $2m\!\int\!\partial_m F\ln(r-2m)\,\dd r$, which by parts is the
horizon dilogarithm $\mathcal{D}_k$ of equation~\eqref{eq:horizon-dilog}. Thus the
advanced/retarded frozen-clock asymmetry isolates the contribution of the
horizon-anchored weight-two sector. For
$m=1,E=1.4,J=2.5$: $\mathcal{A}_\infty=8.7276$, $\mathcal{B}_{\rm hor}=6.1127$ (of which
$1.5301$ is the horizon-log dilogarithm sector), giving
$\delta\varphi/\dot m|_{\rm adv}=14.840$ against
$\delta\varphi/\dot m|_{\rm ret}=2.615$ (both verified to $10^{-15}$); the frozen-clock
difference is then $2\mathcal{A}_\infty=17.455$ and the deviation from antisymmetry
$2\mathcal{B}_{\rm hor}=12.225$ (each $\times|\dot m|$).


\begin{figure}[htbp]
    \centering
    % ALT: A log-log scatter of residual against the adiabatic rate with two fitted straight
    %      lines, one of slope close to one and a steeper one of slope close to two.
    \includegraphics[width=0.72\textwidth]{Immagini/fig_vaidya_offshell.pdf}
    \caption{Vaidya $v$-branch first-order correction, residual against the true
    non-autonomous flow versus $\varepsilon=M\dot m/m$ (log--log). The on-shell term
    alone (dropping the terminally anchored source $S_{\mathcal D}$) is $O(\varepsilon)$;
    the complete extended-Hamiltonian term of equation~\eqref{eq:adiab-exact}, adding the
    off-shell piece, closes the correction to $O(\varepsilon^2)$. Using five rates
    $\varepsilon\in[2.5\times10^{-3},4.0\times10^{-2}]$, the plotted residual is the
    $L^\infty$ orbit-shape norm on $r/M\in[8,11]$. Least-squares fits in
    $(\ln\varepsilon,\ln R)$ give slopes $s_{\rm on}=1.003152$ and
    $s_{\rm full}=1.999997$, with log-fit $L^2$ residuals $2.65\times10^{-3}$ and
    $3.34\times10^{-3}$, respectively.
    Generated by \texttt{VaidyaMetric/vaidya\_first\_order\_offshell.py}.}
    \label{fig:vaidya-offshell}
\end{figure}


% ======================================================================
% Flush pending floats so no figure interrupts the Conclusions.
\FloatBarrier
\subsection{Conjugate points, the field of extremals, and the Maxwell set}
\label{sec:vaidya-jacobi}
Proposition~\ref{prop:hjb} is a criterion; on the frozen family it can be \emph{applied},
because the two obstructions it names---conjugate points and Maxwell points---are both
computable in closed form, and the first turns out never to occur on a monotone arc.

\begin{auxtheorem}[no conjugate point before the turning point]\label{thm:jacobi}
Along the frozen family, with $G=K/\sqrt S$ the shape integrand of
section~\ref{sec:vaidya-adiabatic},
\begin{equation}
   \frac{\partial G}{\partial J}=\frac{r^3(r-2m)^2D_{\Ehat}^{\,2}}{S^{3/2}},
   \label{eq:jacobi-integrand}
\end{equation}
which is strictly positive on the regular domain $r>2m$, $\Ehat>1$, $S>0$. Hence
$\partial\varphi/\partial J>0$ on any radial arc free of turning points: the map
$J\mapsto\varphi(r;J)$ is strictly monotone, the extremals issuing from a common event do
not focus, and no conjugate point occurs there.
\end{auxtheorem}

\begin{proof}
$S=r(r-2m)D_{\Ehat}\,C$ with $C=r^2(r-2m)-J^2D_{\Ehat}$, so
$\partial_JS=-2Jr(r-2m)D_{\Ehat}^{\,2}$ and
\[
   \partial_J\frac{JD_{\Ehat}}{\sqrt S}
   =\frac{D_{\Ehat}}{\sqrt S}+\frac{J^2r(r-2m)D_{\Ehat}^{\,3}}{S^{3/2}}
   =\frac{D_{\Ehat}\big[S+J^2r(r-2m)D_{\Ehat}^{\,2}\big]}{S^{3/2}} .
\]
The bracket collapses: $S+J^2r(r-2m)D_{\Ehat}^{\,2}
=r(r-2m)D_{\Ehat}\big[C+J^2D_{\Ehat}\big]=r^3(r-2m)^2D_{\Ehat}$, which
gives~\eqref{eq:jacobi-integrand}. The numerator is $r^3$ times a square, hence positive,
and $S>0$ on the arc by assumption.
\end{proof}

Combined with strict convexity of the indicatrix (Lemma~\ref{lem:indicatrix}), which makes
the Weierstrass excess function nonnegative, Theorem~\ref{thm:jacobi} embeds each monotone
arc in a Mayer field of extremals: on the region that field covers, the Hilbert invariant
integral constructs the verification function of Proposition~\ref{prop:hjb} implicitly, and
the arc is a \emph{strong} minimizer there. No closed form for $V$ is needed.

\begin{remark}[the Jacobi reduction is already known]\label{rem:jacobi-euler}
Nothing new has to be reduced. The source of~\eqref{eq:jacobi-integrand} is
$N_J=S\,\partial_JK-\tfrac12K\,\partial_JS$, the exact analogue of $N_m$, and the two are
tied by Euler's theorem: $G=K/\sqrt S$ is homogeneous of degree $-1$ in $(r,m,J)$, so
$m\,\partial_mG+J\,\partial_JG=-(rG)'$, i.e.
\begin{equation}
   m\,N_m+J\,N_J=\tfrac12\big[r\,K\,S'-2S\,(rK)'\big].
   \label{eq:jacobi-euler}
\end{equation}
The right-hand side is a pure total derivative, so it contributes no abelian letter, and
the two reductions must satisfy
\begin{equation}
   c^J_k=-\frac{m}{J}\,c^m_k,\qquad
   \mathcal{A}^J=-\frac{rK+m\,\mathcal{A}^m}{J} .
   \label{eq:jacobi-coeffs}
\end{equation}
Every Jacobi coefficient is therefore given by~\eqref{eq:ck-closed}
and~\eqref{eq:Am-closed} already: the conjugate-point analysis and the adiabatic response
run through the same $11\times11$ system, as they must, being the two partial derivatives
of one homogeneous family.
\end{remark}

Theorem~\ref{thm:jacobi} does \emph{not} extend past the turning point, and the reason is
structural rather than technical: on an arc through periapsis the positive
integrand~\eqref{eq:jacobi-integrand} competes with the motion of the turning radius
itself, and the competition is won by neither uniformly.

\begin{auxtheorem}[a caustic of the apsidal family]\label{thm:caustic}
Let $\Phi(J)$ be the azimuth swept by a frozen extremal over one radial excursion from
radius $R$ through periapsis and back to $R$, the turning radius $r_t$ being the root of
$C$ fixed by $J^2=r_t^2(r_t-2m)/D_{\Ehat}(r_t)$. Then $\Phi$ is smooth, strictly positive
on $2m<r_t<R$, and tends to zero at both ends of that range. Consequently $\Phi$ attains an
interior maximum, and at the corresponding $J_c$
\begin{equation}
   \Phi'(J_c)=0 :
   \label{eq:caustic}
\end{equation}
the family develops a \emph{conjugate point} at the return radius. Moreover $\Phi$ is
finite for every admissible $J$, because $C$ has no double root in $r>2m$.
\end{auxtheorem}

\begin{proof}
Write $C(r)=(r-r_t)\,Q(r)$ with
$Q(r)=r^2+(r_t-2m)r+\big[r_t(r_t-2m)-J^2(\Ehat^2-1)\big]$, the division being exact
precisely because $J^2=r_t^2(r_t-2m)/D_{\Ehat}(r_t)$. Substituting $r=r_t+u^2$ turns
$S=(r-r_t)P$ into $u^2P$ with $P(0)=r_t(r_t-2m)D_{\Ehat}(r_t)Q(r_t)\neq0$, so
$G\,\dd r=2K\,\dd u/\sqrt P$ is regular at $u=0$: the apsidal integral is an ordinary
smooth integral in $J$, and differentiation under the integral sign is legitimate.

For the endpoints, $J^2=r_t^2(r_t-2m)/D_{\Ehat}(r_t)\to0$ as $r_t\to2m^+$, and the
integrand carries the overall factor $K=JD_{\Ehat}$, so $\Phi\to0$: a radial plunge sweeps
no azimuth. As $r_t\to R^-$ the integration range $[0,\sqrt{R-r_t}\,]$ shrinks to a point
with bounded integrand, so $\Phi\to0$ again. A positive continuous function vanishing at
both ends of an interval attains an interior maximum, where $\Phi'=0$.

Finiteness is the statement that the turning point is always a \emph{simple} root of $C$.
Eliminating $J^2$ between $C=0$ and $C'=0$ gives $J^2=r(3r-4m)/(\Ehat^2-1)$, positive only
for $r>4m/3$, and the locus $(\Ehat^2-1)r^2-m(\Ehat^2-4)r-4m^2=0$. That quadratic has
positive leading coefficient with value $-4m^2<0$ at $r=0$ and
$\tfrac49(\Ehat^2-1)m^2>0$ at $r=4m/3$, so its only positive root lies in $(0,4m/3)$,
where $J^2<0$. Hence no double root occurs at physical $J$: the integrand behaves as
$(r-r_t)^{-1/2}$ and $\Phi$ converges. This is the same computation that places the
genus degeneration at $r_d<0$ in~\ref{app:degeneration}---the absence of a physical
double root is what both keeps the curve genus two on the arc and keeps the apsidal
integral finite.
\end{proof}

\begin{remark}[the caustic is a simple fold, and it binds first]\label{rem:caustic-num}
Numerically the critical point is a nondegenerate maximum, $\Phi''(J_c)\neq0$, so the
caustic is a simple fold: two extremals coalesce there and none exists beyond. At
$m=1$, $R=10$ one finds $\Ehat=7/5$: $r_t=2.9789$, $J_c=1.3369$, $\Phi_{\max}=0.7432\pi$;
$\Ehat=6/5$: $r_t=2.9387$, $J_c=1.5690$, $\Phi_{\max}=0.7142\pi$; and at $R=20$,
$\Ehat=7/5$: $r_t=3.6876$, $J_c=2.0353$, $\Phi_{\max}=0.8074\pi$. Physically
$\Phi_{\max}$ is the largest azimuth a single in-and-out brachistochrone can span from
radius $R$; beyond it the endpoints must be joined by a multi-excursion extremal.

The critical turning radius scales as $r_t\simeq\sqrt{mR}$, so as $R\to\infty$ the caustic
migrates into the asymptotically flat region and $\Phi_{\max}\to\pi$ from below: the bound
of Theorem~\ref{thm:chord} is sharp, saturated only in the flat limit.
\end{remark}

\begin{remark}[why the sweep bound is exactly what conjugate points need]
\label{rmk:sweep-jacobi}
Theorem~\ref{thm:chord} below is sharper than it looks. In a spherically
symmetric optical geometry the out-of-plane Jacobi field along an equatorial
extremal is generated by the rotation that tilts the orbital plane,
$K=-\sin\varphi\,\partial_\theta-\cot\theta\cos\varphi\,\partial_\varphi$, which
is Killing for the optical metric; a Killing field restricts to a Jacobi field
along any geodesic, and on the equator its optical norm is
$\sqrt{B(r)}\,|\sin\varphi|$ with $B$ the squared area radius. That vanishes at
$\varphi=0$ and next at $\varphi=\pi$, and since the out-of-plane Jacobi fields
vanishing at the launch point form a one-dimensional space, the \emph{first}
out-of-plane conjugate point sits at swept azimuth exactly $\pi$---fixed by the
rotation group, independently of $\Ehat$, of $m$, and of the curvature. Combining
this with Theorem~\ref{thm:jacobi}, which excludes the in-plane conjugate point
on a monotone arc, the two together exclude conjugate points in \emph{all}
directions: the minimality statement is three-dimensional, not merely equatorial.
This is what makes the spherically symmetric sector genuinely stronger than the
rotating one, where no symmetry generates the out-of-plane field.
\end{remark}

\begin{auxtheorem}[the sweep never reaches $\pi$]\label{thm:chord}
For every $m>0$, $\Ehat>1$ and $2m<r_t<R$,
\begin{equation}
   \Phi(r_t)\;\le\;2\arccos\!\big(J/\rho(R)\big)\;<\;\pi,
   \qquad \rho(r)\equiv\sqrt{\mathcal J(r)}=r\sqrt{\frac{r-2m}{D_{\Ehat}(r)}},
   \label{eq:chord-bound}
\end{equation}
with $\mathcal J(r)=r^2(r-2m)/D_{\Ehat}(r)$ the turning function and $J=\rho(r_t)$.
Consequently the axis is never crossed on a single excursion: the caustic of
Theorem~\ref{thm:caustic} is the \emph{only} obstruction there, and the Maxwell mechanism
of Remark~\ref{rem:maxwell} requires more than one radial excursion.
\end{auxtheorem}

\begin{proof}
Eliminating $J^2$ in favour of $\mathcal J$ gives the compact form
$G^2=J^2/\big[r(r-2m)\big(\mathcal J(r)-J^2\big)\big]$. Reparametrizing by
$\rho=\sqrt{\mathcal J}$, and using $r(r-2m)=\mathcal J D_{\Ehat}/r=\rho^2D_{\Ehat}/r$,
\[
   \Phi=4\int_{\rho(r_t)}^{\rho(R)}
   \frac{J\,\sqrt{r/D_{\Ehat}}}{\mathcal J'(r)}\,
   \frac{\dd\rho}{\sqrt{\rho^2-J^2}} ,
\]
so it suffices to prove the pointwise bound $2\rho\sqrt{r/D_{\Ehat}}\le\mathcal J'(r)$,
which reduces $\Phi$ to the Euclidean chord integral
$2\int J\,\dd\rho/(\rho\sqrt{\rho^2-J^2})=2\arccos(J/\rho(R))$.

Writing $b=\Ehat^2-1$, one computes
$\mathcal J'D_{\Ehat}^2=2r\big[br(r-m)+m(3r-4m)\big]$, positive on $r>2m$ since there
$r-m>m$ and $3r-4m>2m$. The bound therefore squares to
\begin{align*}
   &\big[br(r-m)+m(3r-4m)\big]^2-(br+2m)^2\,r(r-2m)\\
   &\qquad=m\big(2br^3+(b^2-6b+5)mr^2+(8b-16)m^2r+16m^3\big)\;\ge\;0 ,
\end{align*}
and substituting $r=2m+s$ with $s>0$ turns the bracket into a polynomial with manifestly
nonnegative coefficients,
\[
   4(1+b)^2+4(1+b)^2\,\tilde s+(b+1)(b+5)\,\tilde s^{2}+2b\,\tilde s^{3},
   \qquad \tilde s=s/m ,
\]
strictly positive for $b>0$, $s>0$. At $m=0$ the difference vanishes identically: the
Euclidean chord is the flat limit, which is where the bound is attained. Finally
$\arccos$ is bounded by $\pi/2$ and $J/\rho(R)>0$ strictly, so
$2\arccos(J/\rho(R))<\pi$.
\end{proof}

\begin{remark}
The comparison must be made in $\rho$, not in $r$: against the Euclidean chord in the
radial variable the pointwise bound \emph{fails}, the rail integrand exceeding the
Euclidean one near periapsis by up to $25\%$. The correct statement is that gravity
strictly shortens the swept azimuth relative to the chord of the \emph{optical} radius
$\rho$, and the deficit is governed by the polynomial above, which vanishes precisely at
$m=0$.
\end{remark}

\begin{remark}[where global optimality does end]\label{rem:maxwell}
The surviving obstruction is the Maxwell set, and the axial symmetry locates it exactly.
Since $S$ is even in $J$, the shape integrand is odd and every clock---proper time
$r^2(r-2m)/\sqrt S$ and advanced time~\eqref{eq:vaidya-clock} alike---is even. Hence
$\varphi(r;-J)=-\varphi(r;J)$ at equal elapsed clock, and the turning radius, a root of
$C$, is shared. The $\pm J$ pair therefore reaches a common point precisely at an axis
crossing $\varphi\equiv0\ (\mathrm{mod}\ \pi)$, and reaches it in equal time: axis
crossings are genuine Maxwell points, where $V$ develops its gradient shock. The closed-form
extremal is thus certified up to its first axis crossing, and beyond it the optimum is
selected by comparing arrival times across the competing windings---the mechanism already
seen in closed FLRW, where the antipodal cut locus of section~\ref{sec:flrw-curvature} is
the same phenomenon in its simplest form.
\end{remark}

\section{Conclusions}
\label{sec:concl}
\paragraph{Proved.} The non-stationary brachistochrone is a well-posed controlled-rail
problem on the regular domain fixed by the two scalar inequalities $W$ timelike and
$\Ehat>|W|$: there the velocity indicatrix is a compact strictly convex oval
(Lemma~\ref{lem:indicatrix}), the invariant $-u\cdot W=\Ehat$ is legitimate via
Pontryagin's principle, with $W$ fixed by the Killing / conformal-Killing / Kodama
hierarchy and the Kodama energy conserved along the Vaidya rail even when free fall does
not. Existence, normality, and the non-autonomous Hamilton--Jacobi verification criterion
hold on that domain. FLRW is the degenerate base at every spatial curvature---freezing
without branch splitting, the extremal being the geodesic of the slice---while in
Vaidya the mass breaks the $t/\tau$ symmetry through its radial gradient.

\paragraph{Conditional and numerical.} The ingoing plunge phenomenology (penetration,
timing, bounce) and the absence of an inversion in the frozen advanced/retarded clock
comparison hold; the negative-rate continuation ($m'<0$ in the ingoing metric) is a formal
continuation, not physical evaporation (which belongs to the outgoing metric). The
first-order adiabatic correction closes in a horizon dilogarithm and is verified against
the true non-autonomous flow to $O(\varepsilon^2)$; its off-shell source has the
universal boundary form $S_{\mathcal D}=[\,rp_r\,]-\lambda$.

\paragraph{Outlook.} The full outgoing-Vaidya boundary-value problem---with its own
Hamiltonians, costates and transversality---remains open, as does a genuine third
(Kodama-observer) clock. The rotating Thakurta--Kerr case, its ergosphere trichotomy,
closed-form separatrices and conformal plunge inversion, is developed in the companion
paper. More broadly, the controlled-rail brachistochrone turns the comparison of
arrival-time and proper-time optimal worldlines into a quantitative probe of how
homogeneous expansion, radial spatial gradients and time-dependent mass flow shape causal
transport in dynamical black-hole spacetimes.


\appendix
\section{Reproducibility and consistency}
\label{app:validation}

\subsection{Reproducing the reductions}
\label{app:reproduce}
\Ithm{4}
\begin{lemma}[polynomial reduction]\label{lem:reduction}
Let the slow derivative of the shape integrand have the form
$\partial_\lambda F=N_\lambda/S^{3/2}$ with $S$ the sextic radicand and $N_\lambda$
polynomial. Then the depth-one reduction
$\partial_\lambda\varphi_0=\mathcal{A}/\sqrt S+\sum_{k=0}^4 c_k U_k$
[equation~\eqref{eq:dm-reduction}] holds if and only if
\begin{equation}
    2\,\mathcal{A}'(r)\,S(r)-\mathcal{A}(r)\,S'(r)
    +2\,S(r)\sum_{k=0}^4 c_k\,r^k \;=\; 2\,N_\lambda(r),
    \qquad \mathcal{A}=\sum_{i=0}^5 a_i r^i,
    \label{eq:reduction-identity}
\end{equation}
a polynomial identity of degree ten. Matching the eleven coefficients of
$r^0,\dots,r^{10}$ gives an $11\times11$ linear system $M$ in the eleven unknowns
$\{a_0,\dots,a_5,c_0,\dots,c_4\}$ whose determinant is, for the Vaidya sextic,
\begin{equation}
    \det M=-32\,(\Ehat^2-1)\,\operatorname{disc}_r S .
    \label{eq:det-disc}
\end{equation}
The system is therefore nonsingular precisely when $S$ is squarefree and $\Ehat\neq1$,
i.e.\ away from the genus-degeneration locus $|J|=\Jc$ and away from the marginal energy;
there its solution is unique and returns the $c_k$ as \emph{rational} functions of the
metric parameters.
\end{lemma}

\begin{proof}
Differentiate the ansatz in $r$ and clear the radical using
$\partial_\lambda F=N_\lambda/S^{3/2}$: the term $\mathcal{A}/\sqrt S$ contributes
$(2\mathcal{A}'S-\mathcal{A}S')/(2S^{3/2})$ and each $U_k=\int r^k/\sqrt S\,\dd r$
contributes $c_k r^k/\sqrt S$, so multiplying through by $2S^{3/2}$
gives~\eqref{eq:reduction-identity}. Both sides are polynomials of degree ten in $r$
(recall $\deg\mathcal{A}=5$, $\deg S=6$), so the identity is equivalent to the vanishing
of eleven coefficients, a linear system in $\{a_i,c_k\}$.

For nonsingularity consider the homogeneous problem
$2\mathcal{A}'S-\mathcal{A}S'+2S\sum_k c_kr^k=0$. It gives $S\mid \mathcal{A}S'$. If $S$
is squarefree then $\gcd(S,S')=1$, hence $S\mid\mathcal{A}$; but
$\deg\mathcal{A}=5<6=\deg S$, forcing $\mathcal{A}=0$ and then $\sum_kc_kr^k=0$, i.e.\
$c_k=0$. The kernel is trivial and the matrix invertible; Cramer's rule exhibits each
$c_k$ as a ratio of determinants in the coefficients of $S$ and $N_\lambda$, hence as a
rational function of the metric parameters.

Conversely, if $S$ is not squarefree---on the degeneration locus, where
$S=(r-r_d)^2Q_4$ with $\deg Q_4=4$---take $\mathcal{A}=(r-r_d)Q_4$, of degree five and
hence admissible. Using $S=(r-r_d)\mathcal{A}$ one computes
$2\mathcal{A}'S-\mathcal{A}S'=(r-r_d)\mathcal{A}\mathcal{A}'-\mathcal{A}^2
=S\,(\mathcal{A}'-Q_4)$, and $\deg(\mathcal{A}'-Q_4)=4$, so
$\big(\mathcal{A},\,\sum_kc_kr^k\big)=\big((r-r_d)Q_4,\,-\tfrac12(\mathcal{A}'-Q_4)\big)$
is a nonzero admissible kernel vector: the system is singular, with a kernel of dimension
one. The determinant formula~\eqref{eq:det-disc} is Lemma~\ref{lem:det-res} below.
\end{proof}

The determinant of the reduction system is not a Vaidya-specific quantity: it is the
resultant of the radicand with its derivative, for any sextic whatsoever.

\begin{auxlemma}[the reduction determinant is a resultant]\label{lem:det-res}
Let $S=\sum_{i=0}^{6}s_ir^i$ be arbitrary and let $M$ be the $11\times11$ matrix, in the
monomial bases, of the linear map
\[
   T:\ \mathbb{P}_5\times\mathbb{P}_4\longrightarrow\mathbb{P}_{10},
   \qquad
   T(\mathcal{A},C)=2\mathcal{A}'S-\mathcal{A}S'+2SC ,
\]
$\mathbb{P}_d$ denoting the polynomials of degree at most $d$. Then
\begin{equation}
   \det M=32\,\mathrm{Res}_r(S,S')=-32\,s_6\,\operatorname{disc}_rS .
   \label{eq:det-res}
\end{equation}
In particular $T$ is invertible if and only if $\deg S=6$ and $S$ is squarefree. For the
Vaidya sextic $s_6=\Ehat^2-1$, which is~\eqref{eq:det-disc}.
\end{auxlemma}

\begin{proof}
Each column of $M$ collects the coefficients of $2(r^i)'S-r^iS'$ or of $2Sr^k$, so every
entry of $M$ is linear in $(s_0,\dots,s_6)$ and $\det M$ is homogeneous of degree eleven in
them. So is $\mathrm{Res}_r(S,S')$: the resultant of a form of degree six with one of
degree five is of degree five in the coefficients of the first and six in those of the
second, and the coefficients of $S'$ are themselves linear in the $s_i$. For degree six the
classical relation reads $\mathrm{Res}_r(S,S')=(-1)^{6\cdot5/2}s_6\operatorname{disc}_rS
=-s_6\operatorname{disc}_rS$.

Next, $\det M$ vanishes on the zero locus of $s_6\operatorname{disc}_rS$. If $s_6=0$ then
$\deg\big(T(\mathcal{A},C)\big)\le9$ for all arguments, so the image lies in the proper
subspace $\mathbb{P}_9$ and $T$ is not surjective. If $s_6\neq0$ and $S$ has a repeated
root $r_d$, write $S=(r-r_d)\mathcal{A}$ with $\mathcal{A}=S/(r-r_d)$ of degree five; since
$r_d$ is at least a double root, $(r-r_d)\mid\mathcal{A}$, say $\mathcal{A}=(r-r_d)Q$ with
$\deg Q=4$. Then
\[
   2\mathcal{A}'S-\mathcal{A}S'
   =\mathcal{A}\big[(r-r_d)\mathcal{A}'-\mathcal{A}\big]
   =(r-r_d)^3\,Q\,Q' ,
\]
so that $T\big(\mathcal{A},\,-\tfrac12(r-r_d)Q'\big)=0$ with $\deg(-\tfrac12(r-r_d)Q')=4$:
a nonzero element of $\ker T$, whence $\det M=0$.

Finally, $\operatorname{disc}_rS$ is irreducible over $\mathbb{Q}$ as a polynomial in the
$s_i$, and it is coprime to $s_6$; both therefore divide $\det M$, hence so does their
product, of degree $1+10=11$. Since $\det M$ is also of degree eleven and is not
identically zero, $\det M=c\,s_6\operatorname{disc}_rS$ for a rational constant $c$, which
a single evaluation fixes: at $S=r^6-1$ one has $\operatorname{disc}_rS=46656$, $s_6=1$
and $\det M=-1492992=-32\cdot46656$, so $c=-32$.
\end{proof}

\begin{remark}[the reduction and the genus are governed by the same discriminant]
Identity~\eqref{eq:det-res} makes Lemma~\ref{lem:reduction} sharp: the depth-one
reduction is uniquely solvable exactly on the genus-two stratum
$\operatorname{disc}_rS\neq0$ and fails exactly where the spectral curve degenerates,
with the rank dropping by exactly one. This is the structural reason why the degeneration
loci need the separate treatment of \ref{app:degeneration}. Independently of the proof
above, the identity was confirmed symbolically in the seven generic coefficients
$(s_0,\dots,s_6)$, and both it and the kernel vector were re-checked on the Vaidya family
in two computer-algebra systems: at $m=1$, $\Ehat=7/5$, $J=5/2$ the determinant is nonzero
with $\gcd(S,S')=1$, while at $|J|=\Jc$ it vanishes with $\operatorname{rank}M=10$ and
$S=(r-r_d)^2Q_4$ exactly.
\end{remark}

The practical consequence is that every coefficient in the adiabatic corrections is fixed
by a \emph{single} polynomial identity that the reader can solve directly; nothing is
fitted. For Vaidya $S$ is the
Schwarzschild-seed sextic and the numerator is
$N_m=S\,\partial_m K-\tfrac12K\,\partial_m S$; the advanced-time clock is entirely
on-curve, its second-kind vector $b$ feeding the $W_{kj}$ of $S$ and its horizon part the
third-kind $\ln\sigma$-quotients.

For the Vaidya mass reduction the linear solve can in fact be carried out once and for
all. Writing $\mathfrak E\equiv\Ehat^{2}$ and introducing the \emph{reduced discriminant}
\begin{equation}
   \Delta\;\equiv\;(\mathfrak E-1)^3J^4+(\mathfrak E^2-20\,\mathfrak E-8)J^2m^2-16m^4,
   \qquad
   \operatorname{disc}_rS=2^{18}\,\mathfrak E^{6}\,J^{10}\,m^{16}\,\Delta,
   \label{eq:reduced-disc}
\end{equation}
the solution of~\eqref{eq:reduction-identity} is
\begin{equation}
\begin{aligned}
   c_0^m&=\frac{4Jm^2\big[3(\mathfrak E-1)J^2+4m^2\big]}{\Delta},&\qquad
   c_4^m&=\frac{2(\mathfrak E-1)(\mathfrak E+8)J}{\Delta},\\[2pt]
   c_1^m&=\frac{4Jm\big[3(\mathfrak E-1)(\mathfrak E-2)J^2+5\mathfrak Em^2\big]}{\Delta},\\[2pt]
   c_2^m&=\frac{J\big[(\mathfrak E-11)(\mathfrak E-1)^2J^2
     +4(\mathfrak E^2-\mathfrak E-15)m^2\big]}{\Delta},\\[2pt]
   c_3^m&=-\frac{J\big[(\mathfrak E-1)^3J^2
     +3(2\mathfrak E^2+9\mathfrak E-20)m^2\big]}{m\,\Delta}.
\end{aligned}
   \label{eq:ck-closed}
\end{equation}
The accompanying algebraic part is likewise explicit,
\begin{equation}
\begin{aligned}
   \mathcal{A}^m&=-\frac{J\,r\,D_{\Ehat}}{m\,\Delta}\;\Pi_{\mathcal A},\\
   \Pi_{\mathcal A}&=(\mathfrak E-1)^3J^4
     +mr\big(\mathfrak E(r-2m)^2+4r(2r-5m)\big)\\
   &\quad+J^2\big((\mathfrak E^2-8\mathfrak E-20)m^2
     +2(\mathfrak E-1)(\mathfrak E-4)mr-(\mathfrak E-1)^2r^2\big),
\end{aligned}
   \label{eq:Am-closed}
\end{equation}
so that~\eqref{eq:dm-reduction} is completely printed: no linear solve is required to
reproduce the Vaidya mass reduction. Substituting~\eqref{eq:ck-closed}
and~\eqref{eq:Am-closed} into~\eqref{eq:reduction-identity} makes it vanish identically in
$(r,m,\Ehat,J)$, which is the statement that they solve it.

Three structural features are visible. Each $c_k^m$ is \emph{odd} in $J$ and depends on the
energy only through $\mathfrak E=\Ehat^2$, as it must, since $N_m$ is odd in $J$ while $S$
is even. The common denominator is the reduced discriminant, so
by~\eqref{eq:det-disc} the reduction degenerates exactly where the curve does. The third
feature is a scaling law, which is not a dimensional observation but a consequence of the
uniqueness established in Lemma~\ref{lem:reduction}.

\begin{auxlemma}[scaling of the reduction]\label{lem:ck-scaling}
Let $\sigma_\alpha$ denote the scaling $(r,m,J)\mapsto(\alpha r,\alpha m,\alpha J)$,
$\alpha>0$, let $S$ be homogeneous of degree six and $N_\lambda$ homogeneous of degree
$\nu$, and assume the hypotheses of Lemma~\ref{lem:reduction}. Then the unique solution
of~\eqref{eq:reduction-identity} is itself homogeneous, of degrees
\begin{equation}
   \deg\mathcal{A}=\nu-5,\qquad \deg c_k=\nu-6-k .
   \label{eq:ck-scaling}
\end{equation}
For the Vaidya mass reduction $K=JD_{\Ehat}$ is homogeneous of degree \emph{two}, hence
$N_m=S\,\partial_mK-\tfrac12K\,\partial_mS$ is homogeneous of degree $\nu=7$, giving
$\deg\mathcal{A}^m=2$ and $\deg c^m_k=1-k$, i.e.\
$c^m_k(\alpha m,\Ehat,\alpha J)=\alpha^{1-k}c^m_k(m,\Ehat,J)$ and equivalently
$c^m_k=m^{1-k}\,C_k(\Ehat^2,J/m)$.
\end{auxlemma}

\begin{proof}
Fix $\alpha>0$ and define
$\tilde{\mathcal{A}}:=\alpha^{5-\nu}\,\mathcal{A}\!\circ\!\sigma_\alpha$ and
$\tilde c_k:=\alpha^{6+k-\nu}\,c_k\!\circ\!\sigma_\alpha$. These are admissible: $\sigma_\alpha$
maps polynomials of degree at most five in $r$ to polynomials of degree at most five, and
the $\tilde c_k$ remain independent of $r$. Evaluating~\eqref{eq:reduction-identity} at
$\sigma_\alpha(r,m,J)$ and using $S\circ\sigma_\alpha=\alpha^6S$,
$(\partial_rS)\circ\sigma_\alpha=\alpha^5\partial_rS$ and
$N_\lambda\circ\sigma_\alpha=\alpha^\nu N_\lambda$ gives
\[
   2\,\alpha^{6}\,(\mathcal{A}'\!\circ\!\sigma_\alpha)\,S
   -\alpha^{5}\,(\mathcal{A}\circ\sigma_\alpha)\,S'
   +2S\sum_k\alpha^{6+k}\,(c_k\circ\sigma_\alpha)\,r^k
   \;=\;2\,\alpha^{\nu}N_\lambda ,
\]
and dividing by $\alpha^{\nu}$ turns the left-hand side into
$2\tilde{\mathcal{A}}'S-\tilde{\mathcal{A}}S'+2S\sum_k\tilde c_kr^k$, because
$\tilde{\mathcal{A}}'=\alpha^{6-\nu}\,\mathcal{A}'\!\circ\!\sigma_\alpha$. Hence
$(\tilde{\mathcal{A}},\tilde c)$ solves the same identity with the same data. By
Lemma~\ref{lem:reduction} that solution is unique, so
$\tilde{\mathcal{A}}=\mathcal{A}$ and $\tilde c_k=c_k$, which is~\eqref{eq:ck-scaling}.
\end{proof}

\begin{remark}
The load-bearing step is uniqueness, not the exponent count: a linear problem whose data
are equivariant under a group action has an equivariant solution only when that solution
is unique. Accordingly~\eqref{eq:ck-scaling} holds exactly on the stratum where
Lemma~\ref{lem:reduction} applies and says nothing on the degeneration locus, where the
solution set is a line and the $c_k$ have poles. The same homogeneity of $S$, used through
Euler's theorem instead, is what fixes $\Jc=m\,j(\Ehat)$ in Lemma~\ref{lem:tracking}.
\end{remark}

Equations~\eqref{eq:ck-closed} and~\eqref{eq:Am-closed} were obtained by solving the
$11\times11$ system symbolically in $(m,\Ehat,J)$; substituted back
into~\eqref{eq:reduction-identity} they annihilate it identically, and they reproduce the
numerical solve below to twelve digits.
A reference solution to check against (polynomial residual $0$ exactly):
\begin{table}[htbp]
    \centering
    \caption{Reference reduction coefficients, obtained by solving
    equation~\eqref{eq:reduction-identity}, at $m{=}1$, $E{=}1.4$, $J{=}2.5$, for the Vaidya
    mass reduction ($\partial_m$). Decimals are shown for readability; each $c_k$ is the
    exact rational solution of a single linear solve (residual $0$ exactly), regenerated
    by \texttt{reproduce\_reductions.py} (exact coefficients in the reproducibility
    package).}
    \label{tab:reduction-ck}
    \small
    \begin{tabular}{lrrrrr}
        \hline\hline
        reduction & $c_0$ & $c_1$ & $c_2$ & $c_3$ & $c_4$\\
        \hline
        Vaidya, $\partial_m$ & $-0.87153$ & $-0.35970$ & $1.03538$ & $0.21292$ & $-0.18939$\\
        \hline\hline
    \end{tabular}
\end{table}

Given the $c_k$, the weight-two assembly is mechanical: the clock supplies the vector
$b$ (proper-time $T_\tau=U_3-2m\,U_2$, i.e.\ $b=(0,0,-2m,1,0)$), the algebraic
weights are $Q_{kj}=c_kb_j-c_jb_k$, the weight-two
letters are $W_{kj}=\int(U_k\dd U_j-U_j\dd U_k)$, and the elementary remainders
$\mathcal{I}_j=\int\mathcal{A}r^j/S\,\dd r$ integrate to logarithms. For Vaidya the
same steps with the advanced-time clock give the fully assembled
form~\eqref{eq:vaidya-full}; the horizon dilogarithm closes by one integration by
parts, equation~\eqref{eq:horizon-dilog}. Table~\ref{tab:script-map} maps each Vaidya result to
the script in the repository~\cite{Rosignoli2026CodeJMP} that generates and verifies it.
\begin{table}[htbp]
    \centering
    \caption{Result-to-script map with the achieved numerical residual, so each closed
    representation can be regenerated and re-checked independently. Scripts are under
    \texttt{VaidyaMetric/} in the archived repository~\cite{Rosignoli2026CodeJMP}.}
    \label{tab:script-map}
    \small
    \begin{tabular}{@{}p{0.42\textwidth}>{\raggedright\arraybackslash}p{0.34\textwidth}r@{}}
        \hline\hline
        result & script & residual\\
        \hline
        Vaidya reduction $c_k^m$
            & \texttt{vaidya\_dMF\_reduction.py} & $10^{-15}$\\
        Vaidya clock $v$; assembled $\delta\varphi$~\eqref{eq:vaidya-full}
            & \texttt{vaidya\_fully\_explicit.py} & $2\times10^{-13}$\\
        Vaidya $\tau$-branch assembly
            & \texttt{vaidya\_tau\_assembly.py} & $10^{-14}$\\
        horizon dilogarithm $\mathcal{D}_k=U_k\ln-\mathcal{G}_k$
            & \texttt{vaidya\_horizon\_dilog.py} & $10^{-12}$\\
        genus-degeneration elliptic form $U_0(\sigma,\zeta)$; $\tau$
            & \texttt{vaidya\_ell\_dilog\_match.py} & $10^{-14}$\\
        advanced/retarded frozen-clock decomposition $\mathcal{A}_\infty,\mathcal{B}_{\rm hor}$
            & \texttt{vaidya\_asymmetry.py} & $10^{-15}$\\
        complete first-order correction~\eqref{eq:adiab-exact} (on-shell $+$ off-shell)
            & \texttt{vaidya\_first\_\allowbreak order\_\allowbreak offshell.py} & slope $2.00$\\
        \hline\hline
    \end{tabular}
    \\[2pt]
    {\footnotesize SHA-256 (first 12 hex) of
    \texttt{vaidya\_first\_order\_offshell.py}: \texttt{5dc9fe177f2f}.}
\end{table}

\subsection{Consistency residuals}
\label{app:residuals}
The closed forms and reductions above were checked symbolically and numerically. The
Kodama energy is held to $|{-}u\cdot K-\Ehat|<7\times10^{-16}$ along the Vaidya rail; the
stationary limit $m'=0$ reproduces the closed-form Schwarzschild ($a{=}0$ Kerr) orbit to
$10^{-12}$ (figure~\ref{fig:vaidya-kerra0}); each branch is a local minimizer of its own
travel time within the tested perturbation family, $t$ and $\tau$ being distinct curves
(figure~\ref{fig:verifmin}); and the first-order
adiabatic correction closes to $O(\varepsilon^2)$ against the true non-autonomous flow
(figure~\ref{fig:vaidya-offshell}).










\begin{figure}[tbp]
    \centering
    % ALT: A numerically integrated orbit and a closed-form reference orbit plotted together
    %      and visually indistinguishable, with their difference staying below 1e-12.
    \includegraphics[width=\textwidth,height=0.68\textheight,keepaspectratio]{Immagini/fig_vaidya_kerr_a0.pdf}
    \caption{Stationary-limit check ($m'=0$): the numerical orbit matches the
    closed-form Kerr $a=0$ solution to a sup-norm residual $<10^{-12}$.}
    \label{fig:vaidya-kerra0}
\end{figure}


\begin{figure}[tbp]
    \centering
    % ALT: Travel time against perturbation amplitude, one smooth curve per clock branch,
    %      each with a single minimum at zero amplitude and the two curves separated.
    \includegraphics[width=\textwidth,height=0.68\textheight,keepaspectratio]{Immagini/fig_verifica_minimo_brachi.pdf}
    \caption{Local perturbation check: within the tested family each branch is a minimizer
    of its own travel time; $t$ and $\tau$ are distinct curves.}
    \label{fig:verifmin}
\end{figure}









\section{Explicit adiabatic forms on the genus-degeneration loci}
\label{app:degeneration}
This appendix collects in one place the fully explicit first-order adiabatic
brachistochrone on the loci where the branch sextic acquires a double root and drops to
genus one. For \emph{spherical} Vaidya the relevant value $J=\Jc$ is the \emph{algebraic}
genus-degeneration $J_{\rm deg}$---the double root sits at $r_d<0$
(a negative-radius double root), so it is \emph{not} a physical, accessible separatrix,
though the genus-two$\,\to\,$genus-one reduction is valid there all the same; we give both
the proper-time ($\tau$) and advanced-time ($v$) branches with every coefficient in closed
form. All objects live on
the elliptic curve $E:\,w^2=Q_4(r)$ into which the branch sextic collapses. The adiabatic expansion is
$\varphi(r)=\varphi_0(r)+\dot\lambda\,\delta\varphi(r)+O(\dot\lambda^2)$, $\lambda=m$.

\paragraph{Curve and special functions.}
$\Jc=5\sqrt{3011/3072+581\sqrt{249}/9216}$ (at $E=7/5,\,m=1$; in general $\Jc=m\,j(E)$),
$S=(r-r_d)^2Q_4(r)$ with double root $r_d$ and quartic $Q_4$ (roots
$e_1<e_2<e_3=2m<e_4$, leading coefficient $a_4$). The Abel coordinate is
$z(r)=\int_{r_0}^r\dd r'/\sqrt{Q_4}$; $\sigma,\zeta,\wp$ are the Weierstrass functions
of $E$ (periods $2\omega_1,2i\omega_{\rm im}$, built from $\theta_1$). Marked points:
$\pm z_d$ (images of $r_d$), $\pm z_\infty$ (of $r=\infty$), the two-torsion turning
image $z=0$ ($r=e_4$) and horizon image $z=i\omega_{\rm im}$ ($r=e_3=2m$).

\paragraph{Zeroth order (frozen orbit), verified to $10^{-12}$.}
\begin{equation}
\varphi_0(r)=\Jc\big[(E^2-1)z+((E^2-1)r_d+2m)\,U_0(z)\big],\;\;
U_0=\rho[\ln\sigma(z{-}z_d)-\ln\sigma(z{+}z_d)]+C_{U_0}z,
\end{equation}
$\rho=1/\sqrt{Q_4(r_d)}$, $U_k=\int_{r_0}^r r^k/\sqrt S\,\dd r$. Both the shape of $U_0$
and the constant $\rho$ are forced, and the higher $U_k$ reduce to it.

\begin{auxlemma}[third-kind letter and Hermite reduction]\label{lem:U0}
On the degenerate curve $S=(r-r_d)^2Q_4$, with Abel coordinate
$z(r)=\int_{r_0}^r\dd r'/\sqrt{Q_4}$ and $\pm z_d$ the preimages of $r_d$:
\begin{enumerate}
\item[(i)] $\displaystyle U_0(z)=\rho\big[\ln\sigma(z-z_d)-\ln\sigma(z+z_d)\big]+C_{U_0}z$
with $\rho=1/\sqrt{Q_4(r_d)}$;
\item[(ii)] for $k\ge1$, $\;U_k=V_k+r_d^{\,k}\,U_0$, where
$V_k=\int P_{k-1}(r(z))\,\dd z$ with $P_{k-1}$ the quotient of $r^k$ by $r-r_d$, an
elliptic integral of the second kind reducible to $\zeta,\wp$ and their derivatives at
$z_\infty$ (Hermite).
\end{enumerate}
\end{auxlemma}

\begin{proof}
Since $\sqrt S=(r-r_d)\sqrt{Q_4}$ and $\dd r=\sqrt{Q_4}\,\dd z$, every letter collapses to
$U_k=\int r^k\,\dd z/(r-r_d)$. Dividing, $r^k/(r-r_d)=P_{k-1}(r)+r_d^{\,k}/(r-r_d)$, which
is (ii), with (i) the remaining case $k=0$.

For (i), differentiate: $U_0'(z)=1/(r(z)-r_d)$. Now $r(z)$ is elliptic of order two, so
$U_0'$ is elliptic with simple poles exactly at the two preimages $\pm z_d$ of $r_d$; and
because $r'(z)=\sqrt{Q_4(r)}$ by construction of $z$, the residues are
\[
   \operatorname*{Res}_{z=\pm z_d}\frac{1}{r(z)-r_d}=\frac{1}{r'(\pm z_d)}
   =\pm\frac{1}{\sqrt{Q_4(r_d)}}=\pm\rho .
\]
The candidate $\rho\big[\zeta(z-z_d)-\zeta(z+z_d)\big]$ is elliptic---the quasi-period
$2\eta$ cancels in the difference---and has simple poles at the same two points with the
same residues $\pm\rho$. Their difference is therefore an elliptic function with no poles,
hence a constant, which is $C_{U_0}$. Integrating and using $\zeta=(\ln\sigma)'$ gives
(i); the constant $\rho$ is not fitted but is the residue itself.
\end{proof}

\begin{remark}
Part (ii) is the degeneration statement quoted in \ref{app:firstorder}, $U_k\to V_k+
r_d^{\,k}\Pi_{r_d}$, with $\Pi_{r_d}=U_0$ the third-kind letter: the whole tower of
explicit coefficients below rests on the single residue computation above. It also
explains the shape of $r(z)=c_r-a_4^{-1/2}[\zeta(z-z_\infty)-\zeta(z+z_\infty)]$: an
order-two elliptic function is fixed by its two simple poles---here the preimages
$\pm z_\infty$ of $r=\infty$---and the same difference-of-$\zeta$ construction applies.
\end{remark}

\paragraph{First order: source and antiderivative.}
$\partial_m F=N_m/S^{3/2}$, $N_m=S\,\partial_mK-\tfrac12K\,\partial_mS$, $K=\Jc\,D_E$,
$D_E=(E^2-1)r+2m$; on the curve $\partial_m F\,\dd r=R(z)\,\dd z$ with
$R=N_m(r(z))/[(r(z)-r_d)^3Q_4(r(z))]$ and $r(z)=c_r-a_4^{-1/2}[\zeta(z-z_\infty)-
\zeta(z+z_\infty)]$. Its antiderivative is equation~\eqref{eq:vaidya-sep-G},
$\tilde G(z)=C_0 z+\sum_{a}[b_1^a\ln\sigma(z-a)-b_2^a\zeta(z-a)-\tfrac12 b_3^a
\wp(z-a)]$, $a\in\{\pm z_d,0,i\omega_{\rm im}\}$, with closed-form residues: at the
triple pole $z_d$, writing $s_d=\sqrt{Q_4(r_d)}$, $F=N_m/Q_4$, $a_1=Q_4'(r_d)/(4s_d)$,
$a_2=Q_4''(r_d)/12$, $h_0=F(r_d)$, $h_1=F'(r_d)s_d$,
$h_2=\tfrac12F''(r_d)s_d^2+\tfrac14F'(r_d)Q_4'(r_d)$,
\begin{equation}
\begin{aligned}
b_3^{z_d}&=\frac{h_0}{s_d^3},\qquad
b_2^{z_d}=\frac{h_1-3a_1h_0}{s_d^3},\qquad
b_1^{z_d}=\frac{h_2-3a_1h_1+(6a_1^2-3a_2)h_0}{s_d^3},\\
b_2^{e_i}&=\frac{4N_m(e_i)}{(e_i-r_d)^3Q_4'(e_i)^2},
\end{aligned}
\end{equation}
($e_i=e_4,e_3$ the two-torsion poles) and the constant closing at $z_\infty$, where
$R$ vanishes: $C_0=-\sum_a[b_1^a\zeta(z_\infty-a)+b_2^a\wp(z_\infty-a)-\tfrac12
b_3^a\wp'(z_\infty-a)]$. None of these coefficients is fitted; each is a Laurent
coefficient, as the following makes explicit.

\begin{auxlemma}[residues of the source antiderivative]\label{lem:residues}
With $\varepsilon=z-z_d$ and $s_d=\sqrt{Q_4(r_d)}$:
\begin{enumerate}
\item[(i)] $r-r_d=s_d\,\varepsilon\big(1+a_1\varepsilon+a_2\varepsilon^2+O(\varepsilon^3)\big)$
with $a_1=Q_4'(r_d)/(4s_d)$ and $a_2=Q_4''(r_d)/12$;
\item[(ii)] the Laurent coefficients of $R$ at $z_d$ are the $b_3^{z_d},b_2^{z_d},b_1^{z_d}$
displayed above;
\item[(iii)] at a simple root $e$ of $Q_4$ the function $r(z)$ is \emph{even} about the
corresponding two-torsion point, so $R$ has there a double pole and \emph{no} residue,
with coefficient $b_2^{e}=4N_m(e)/[(e-r_d)^3Q_4'(e)^2]$;
\item[(iv)] $\tilde G'=R$ then holds with the displayed $C_0$.
\end{enumerate}
\end{auxlemma}

\begin{proof}
(i) Insert the ansatz into $(\dd r/\dd z)^2=Q_4(r)$, which defines $z$, and match orders in
$\varepsilon$; the first two corrections come out as stated.

(ii) Expanding $F=N_m/Q_4$ about $r_d$ and using (i),
$F=h_0+h_1\varepsilon+h_2\varepsilon^2+O(\varepsilon^3)$ with exactly the $h_j$ displayed
above---note $h_2=\tfrac12F''s_d^2+F's_da_1$ and $a_1s_d=Q_4'/4$, which is the second form
given. Meanwhile
$(r-r_d)^{-3}=(s_d^3\varepsilon^3)^{-1}\big[1-3a_1\varepsilon+(6a_1^2-3a_2)\varepsilon^2
+O(\varepsilon^3)\big]$. Multiplying the two series and reading off the coefficients of
$\varepsilon^{-3},\varepsilon^{-2},\varepsilon^{-1}$ gives $b_3^{z_d},b_2^{z_d},b_1^{z_d}$.

(iii) Near a simple root, $\delta=r-e$ obeys $(\dd\delta/\dd z)^2=Q_4'(e)\delta+O(\delta^2)$.
Matching orders gives $\delta=Q_4'(e)\varepsilon^2/4+O(\varepsilon^4)$ with every
\emph{odd} coefficient vanishing---at order $\varepsilon^3$ the matching reads
$3Q_4'(e)k_3=Q_4'(e)k_3$, forcing $k_3=0$, and likewise at each odd order. Hence $r$, and
therefore $R$, is even about that point: only even negative powers occur, so there is a
double pole and no residue, and its coefficient follows from
$Q_4(r)\simeq Q_4'(e)^2\varepsilon^2/4$.

(iv) Differentiating term by term, $\dd(\ln\sigma)/\dd z=\zeta$, $\zeta'=-\wp$ and
$\wp'(u)=-2u^{-3}+O(u)$, so the three families contribute principal parts
$b_1^a(z-a)^{-1}$, $b_2^a(z-a)^{-2}$ and $b_3^a(z-a)^{-3}$ respectively: $\tilde G'$ and
$R$ have identical principal parts at every pole, so their difference is a pole-free
elliptic function, i.e.\ a constant. Evaluating at $z_\infty$, where $R$ vanishes, fixes it
to the displayed $C_0$.
\end{proof}

\paragraph{First order, proper-time branch (verified to $10^{-9}$).}
Clock $T_\tau=U_3-2mU_2$; block form
\begin{equation}
\delta\varphi_\tau=\tfrac12\,\tilde G(z)\,T_\tau(r)
+\tfrac12\sum_{i<j}(c_i d_j-c_j d_i)\,A[e_i,e_j],
\end{equation}
$A[e_i,e_j]=\int(e_i\,\mathrm{Pe}_j-e_j\,\mathrm{Pe}_i)\dd z$ the antisymmetric
elliptic dilogarithms (Brown--Levin $\tilde\Gamma$), over the eight even blocks
$\{1,Z_{z_d},P_{z_d},P'_{z_d},\wp(z),\wp(z{-}i\omega_{\rm im}),P_{z_\infty},
Z_{z_\infty}\}$ with $Z_a=\zeta(z{-}a)-\zeta(z{+}a)$,
$P_a=\wp(z{-}a)+\wp(z{+}a)$, $P'_a=\wp'(z{-}a)-\wp'(z{+}a)$. Coefficients: $c_i$ (the
$\tilde G$ residues), $d_j$ (clock residues $e_1^{z_d}=(r_d^3-2m\,r_d^2)/s_d$,
$e_2^{z_\infty}=1/a_4$, $e_1^{z_\infty}=-(2B+r_d-2m)/\sqrt{a_4}$, where
$B=-a_3/(4a_4)$ is the constant term of $r(z)$ at $z_\infty$ and $a_3$ the cubic
coefficient of $Q_4$) and the constant
$C_e=(e_4^3-2m\,e_4^2)/(e_4-r_d)+2e_1^{z_d}\zeta(z_d)-2e_2^{z_\infty}\wp(z_\infty)
+2e_1^{z_\infty}\zeta(z_\infty)$---all closed form.

\begin{auxlemma}[block form of the first-order correction]\label{lem:blocks}
Let $\tilde G$ be the antiderivative of $R$ constructed in
Lemma~\ref{lem:residues}, let $T$ be any clock, and expand both in the even block
basis, $\tilde G=\sum_ic_i\,\mathrm{Pe}_i$ and $T=\sum_jd_j\,\mathrm{Pe}_j$, with
$\mathrm{Pe}$ the primitive of a block. Then
\begin{equation}
   \delta\varphi=\int R\,T\,\dd z
   =\tfrac12\,\tilde G\,T+\tfrac12\sum_{i<j}(c_id_j-c_jd_i)\,A[e_i,e_j],
   \qquad A[e_i,e_j]=\int(e_i\mathrm{Pe}_j-e_j\mathrm{Pe}_i)\,\dd z .
   \label{eq:block-form}
\end{equation}
For $T=T_\tau=U_3-2mU_2$ the coefficients are those listed above; in particular
$e_1^{z_d}=(r_d^3-2m\,r_d^2)/s_d$, $e_2^{z_\infty}=1/a_4$ and
$e_1^{z_\infty}=-(2B+r_d-2m)/\sqrt{a_4}$ with $B=-a_3/(4a_4)$.
\end{auxlemma}

\begin{proof}
Since $\tilde G'=R$, integrating by parts twice and averaging gives the elementary
identity $\int\tilde G'T=\tfrac12\tilde GT+\tfrac12\int(\tilde G'T-\tilde GT')$, which is
the displayed boundary term plus a manifestly antisymmetric integrand. Substituting the
two expansions,
$\tilde G'T-\tilde GT'=\sum_{i,j}c_id_j\,(e_i\mathrm{Pe}_j-\mathrm{Pe}_ie_j)$; the
diagonal terms cancel and the transposed pair combines to $(c_id_j-c_jd_i)$, giving
\eqref{eq:block-form}.

For the clock coefficients, Lemma~\ref{lem:U0}(ii) gives
$T_\tau=(V_3-2mV_2)+(r_d^3-2m\,r_d^2)\,U_0$, and $U_0$ carries residue $\rho=1/s_d$ at
$z_d$, whence $e_1^{z_d}=(r_d^3-2m\,r_d^2)/s_d$. At $z_\infty$ expand the clock integrand
$(r^3-2mr^2)/(r-r_d)$ using the local behaviour of $r(z)$: matching orders in
$(\dd r/\dd z)^2=Q_4(r)$ gives
\[
   r(z)=-\frac{1}{\sqrt{a_4}\,(z-z_\infty)}+B+O(z-z_\infty),
   \qquad B=-\frac{a_3}{4a_4},
\]
and reading off the double-pole coefficient and the residue returns $e_2^{z_\infty}=1/a_4$
and $e_1^{z_\infty}=-(2B+r_d-2m)/\sqrt{a_4}$. The constant $C_e$ is then fixed exactly as
$C_0$ was in Lemma~\ref{lem:residues}.
\end{proof}

\paragraph{First order, advanced-time branch (verified to $10^{-8}$).}
Clock $v=E\,U_3+r+2m\ln(r-2m)$;
$\delta\varphi_v=\tilde G(z)\,v-\int_{z_0}^z\tilde G\,v_z\,\dd z$ with
$v_z=E\,r^3/(r-r_d)+r\sqrt{Q_4}/(r-2m)$. The first term reproduces the proper-time
dilogarithms; the second is odd on $E$ with a simple pole at $z=i\omega_{\rm im}$ and
supplies the horizon dilogarithm.

\subsection{Adiabatic derivative of the degeneration family}
\label{app:tracking}
As the slow parameter drifts, the genus-degeneration value itself moves,
$J=\Jc(\lambda)$. This is \emph{not} the evolution of a single torque-free extremal:
the axial costate $J$ is conserved along every extremal
(section~\ref{sec:principle}), so one trajectory cannot follow the moving locus by
changing $J$ unless an external azimuthal torque---a modified control law with
$\dd J/\dd\lambda\neq0$---is supplied. The object below is instead the total
$\lambda$-derivative of the \emph{one-parameter family of marginal (degeneration)
solutions} along $J=\Jc(\lambda)$: the adiabatic derivative of the family, which
packages how the closed form deforms as the locus drifts.

\Ithm{6}
\begin{lemma}[degeneration-family derivative]\label{lem:tracking}
Let $r_d(\lambda)$ be the double root of $S$ along the degeneration locus
$J=\Jc(\lambda)$, and set
\begin{equation}
N_{\rm tot}=N+\frac{\dd\Jc}{\dd\lambda}\,N_J,\qquad
N_J=S\,\partial_JK-\tfrac12K\,\partial_JS.
\label{eq:Ntot-track}
\end{equation}
Then the block form and residue formulae of \ref{app:degeneration} hold verbatim under
$N\to N_{\rm tot}$, with $\varphi_0$, the clocks and the dilogarithm structure unchanged;
and the tracked source vanishes at the double root,
\begin{equation}
N_{\rm tot}(r_d)=-\tfrac12K(r_d)\Big[\partial_\lambda S
+\tfrac{\dd\Jc}{\dd\lambda}\,\partial_JS\Big](r_d)
=\tfrac12K(r_d)\,S'(r_d)\,\frac{\dd r_d}{\dd\lambda}=0,
\label{eq:track-cancel}
\end{equation}
so that the triple-pole residue drops, $b_3^{\rm track}=0$.
\end{lemma}

\begin{proof}
Evaluating~\eqref{eq:Ntot-track} at $r=r_d$, the terms proportional to $S$ drop because
$S(r_d)=0$, leaving the first equality in~\eqref{eq:track-cancel}. For the second, note
that the Vaidya sextic is \emph{jointly homogeneous of degree six} in $(r,m,J)$,
\[
   S(\alpha r,\alpha m,\alpha J)=\alpha^6\,S(r,m,J),
   \qquad\text{hence}\qquad
   r\,\partial_rS+m\,\partial_mS+J\,\partial_JS=6\,S
\]
by Euler's theorem. Homogeneity also fixes the scaling of the degeneration locus,
$\Jc=m\,j(\Ehat)$ and $r_d=m\,\rho(\Ehat)$, so that $\dd\Jc/\dd m=\Jc/m$ and
$\dd r_d/\dd m=r_d/m$. Evaluating Euler's identity at the double root, where both $S$ and
$\partial_rS=S'$ vanish, leaves $m\,\partial_mS+J\,\partial_JS=0$, that is
\[
   \Big[\partial_mS+\frac{\Jc}{m}\,\partial_JS\Big](r_d)
   =\Big[\partial_\lambda S+\frac{\dd\Jc}{\dd\lambda}\,\partial_JS\Big](r_d)=0 ,
\]
which is the bracket in~\eqref{eq:track-cancel}; equivalently it equals
$-S'(r_d)\,\dd r_d/\dd\lambda$, obtained by differentiating the double-root condition
$S(r_d(\lambda),\Jc(\lambda);\lambda)=0$ totally in $\lambda$, and this vanishes because
$S'(r_d)=0$. Hence $N_{\rm tot}(r_d)=0$, and since the triple-pole residue is proportional
to $N_{\rm tot}(r_d)$ it vanishes as well.
\end{proof}

\begin{remark}[independent verification]
Euler's identity for $S$ was confirmed to reduce to zero identically in
\textsc{Mathematica} and \textsc{SymPy}. At the concrete locus $m=1$, $\Ehat=7/5$---where
$\Jc=5\sqrt{3011/3072+581\sqrt{249}/9216}=7.0266237403890\ldots$, reproducing the value
quoted in \ref{app:degeneration}, and $r_d=-3.3637111847170\ldots$---one finds
$N(r_d)=1.16239\times10^{4}$ and $N_J(r_d)=-1.65427\times10^{3}$ separately nonzero, while
$N_{\rm tot}(r_d)=0$ exactly in both systems.
\end{remark}

Following the moving genus-degeneration locus therefore softens the strongest
singularity, precisely where the double root travels. The cancellation is verified
symbolically in Mathematica, $N_{\rm tot}(r_d)\equiv0\pmod{S,\,S'}$. The two Vaidya
branches (proper-time $\tau$ and advanced-time $v$) differ only in $\dd\Jc/\dd\lambda$
and $N_J$.

\paragraph{Vaidya ($\lambda=m$).}
The linear scaling $\Jc=m\,j(E)$ gives $\dd\Jc/\dd m=\Jc/m$, with
$N_J=S\,D_E+\Jc^2\,r(r-2m)D_E^3$ (verified to $10^{-8}$).

\paragraph{Fully symbolic coefficients, generic-genus $q$-series, and independent checks.}
Away from the degeneration locus (generic $J$, genus two) every block coefficient is a
\emph{rational function of $(m,E,J)$}: the source coefficients $c_k$,
$Q_{kj}=c_kb_j-c_jb_k$, and the canonical pole coefficients $g_i$ (Taylor
coefficients of $q_6^{-1/2}$ at the odd-model infinity), with
$g_0=1/\sqrt{E^2-1}$ and $g_1=(2E^2-3)/(E^2-1)^{3/2}$ set by the leading behaviour
($E$-only), while $g_2$ and the $c_k$ carry the full $J$ dependence. The
algebraic remainder is elementary in closed form,
$T_{\rm alg}=P(r)+\sum_i\mathrm{res}_i\ln(r-r_i)$ over the six roots of $S$
(three residues vanish, at $r=0,\,2m,\,2m/(1-E^2)$). The genus-two dilogarithm is a
convergent \emph{nome ($q$-)series}: after Siegel reduction $\tau'=-\tau^{-1}$ (all
nomes ${<}1$), $\Lambda=\int Q\,\dd\mathcal A+\int L\,\dd\mathcal A$ with $Q$
elementary and $L=\log[\theta[\delta](w-e_+)/\theta[\delta](w-e_-)]$ a genus-two
Kronecker--Eisenstein series that converges geometrically (order $N{=}4$ already
matches to $10^{-13}$) and does \emph{not} reduce to classical dilogarithms (there is
no genus-two Jacobi product formula). All these symbolic statements are
\emph{independently reproduced in Mathematica} (native $\wp$, Gröbner reduction,
Laurent extraction): the degeneration residues to $10^{-16}$--$10^{-38}$, the generic-$J$
reductions to an exact polynomial identity, and the tracking
theorem~\eqref{eq:track-cancel}.

\paragraph{Transcendence.}
None of these is elementary. The $U_k$ are third-kind elliptic integrals, $\tilde G$
lives in $\sigma,\zeta,\wp$, and the weight-two content is the elliptic dilogarithm,
conjecturally irreducible to weight one (theta-divisor obstruction $\Theta=W_{g-1}$). The
advanced-time horizon logarithm is elementary only as a $\dd\log$ letter; integrated
against the elliptic $\partial_m F$ it yields the horizon elliptic dilogarithm, of the
same transcendental weight as the proper-time case---the horizon merely relocates the
dilogarithm's anchor from infinity to the two-torsion point $r=2m$. The sole
numerical ingredient anywhere is the optional compression of the weight-two content to
an apparent minimal five-dimensional dilogarithm basis (a numerically stable rank five,
singular values falling to $10^{-79}$---strong evidence, \emph{not} a proof, and
curve-specific), which the block forms above do not require.

\FloatBarrier


\section{Complete first-order extended-Hamiltonian correction}
\label{app:firstorder}

This appendix proves Theorem~\ref{thm:firstorder}, stated in section~\ref{sec:vaidya-adiabatic}: it derives
equation~\eqref{eq:adiab-exact}---the complete first-order (on-shell \emph{plus} off-shell)
adiabatic correction to the orbit shape---from the extended Hamiltonian in a
metric-independent way, and closes it explicitly for the Vaidya branches so that a reader
can reconstruct it without the accompanying code. Here $\Theta$ is the Euler operator that
generates the drift of the frozen family through the slow parameter, and the drift
generator $\mathcal D$ equals $\Theta$ \emph{unless} the slow parameter rescales the
canonical momenta, in which case it acquires a momentum-normalization term. For
single-parameter Vaidya the mass is a metric modulus (no rescaling), so
$\mathcal D=\Theta=m\partial_m$ and the source collapses to the universal boundary form
$S_{\mathcal D}=[\,rp_r\,]-\lambda$ derived below. If a slow parameter also rescales a
canonical momentum, the drift generator requires the corresponding dilation term; that
different deformation class is not used in the Vaidya calculation below.

\paragraph{Notation and the slow datum.} $\Ehat$ and $J=p_\varphi$ are the
\emph{conserved} control charges (Kodama energy and axial costate). For Vaidya the
background coordinate is the advanced time $v$. Before expanding, choose a fixed
reference mass $M$ and introduce the dimensionless slow coordinate $s=v/M$; radii and
masses in this appendix are likewise measured in units of $M$. We then write
$m(s)=m_0(1+\varepsilon s+\cdots)$ with the unambiguous dimensionless rate
\[
    \varepsilon\equiv\frac{1}{m}\frac{\dd m}{\dd s}
    =\frac{M\dot m}{m},\qquad \dot m\equiv\frac{\dd m}{\dd v}.
\]
The operator generating the drift of the frozen family is the
Euler operator
\begin{equation}
    \Theta\equiv m\,\partial_m\quad(\text{Vaidya}),
    \label{eq:theta-op}
\end{equation}
since $\dd m/\dd s=+\varepsilon m$ at leading order.

\paragraph{Frozen data.} At fixed $s$ the frozen Hamiltonian $H_0(r,p_r;\lambda)=0$
fixes: the \emph{ingoing} frozen momentum $p_{r0}(r)$, the root of $H_0=0$ with
$H_{0,p_r}<0$ (the outgoing arc is its mirror across the turning point); the shape
integrand $G_0(r)=\dd\varphi/\dd r=H_{0,J}/H_{0,p_r}$; and the clock density
$\dd\lambda/\dd r=1/H_{0,p_r}$, with $\lambda$ oriented \emph{inward from the
free-arrival endpoint}.

\paragraph{Drift generator.} Adjoining the slow clock $s$ with conjugate $p_s$, the
extended Hamiltonian $\mathcal H=p_s+H(r,p_r;\lambda(s))=0$ generates the \emph{canonical}
flow $\dd r/\dd s=H_{p_r}$, $\dd p_r/\dd s=-H_r$, and the shell-departure rate is
$\dd H/\dd s=-\varepsilon\,\mathcal D H$ with $\mathcal D$ the Euler operator generating
the frozen-family drift. When the slow parameter enters only as a metric modulus
(Vaidya) it does not rescale the momenta, so $\mathcal D=\Theta=m\partial_m$. When
instead a slow parameter rescales a normalized momentum $P_r$, the rescaling is
time-dependent (conformally symplectic), \emph{not} canonical: the chain rule adds a
dilation term and promotes $\Theta$ to a full Euler operator containing
$P_r\partial_{P_r}$. Here no such rescaling occurs: the canonical pair is $(r,p_r)$ and
$\mathcal D=\Theta$.

\paragraph{Extended system and first order.} Adjoining $s$ with conjugate $p_s$,
$\mathcal H=p_s+H(r,p_r;\lambda(s))=0$; free-arrival transversality sets $H=0$ (hence
$p_s=0$) at $\lambda=0$. Along the flow $\dd p_s/\dd\lambda=\varepsilon\,\mathcal D H$, so
the \emph{terminally-anchored} source is
\begin{equation}
    p_s^{(1)}(\lambda)\equiv S_{\mathcal D}(\lambda)
    =\int_0^{\lambda}\mathcal D H\,\dd\lambda',
    \label{eq:pS-app}
\end{equation}
and, linearizing $H$ at fixed $r$, the momentum-shell shift
\begin{equation}
    \delta p_r=\frac{\lambda\,\Theta H-S_{\mathcal D}}{H_{p_r}},
    \label{eq:dpr-app}
\end{equation}
the axial costate remaining conserved ($\delta J=0$). The shape varies as
$\dd\varphi/\dd r=G_0+\varepsilon[\partial_{p_r}G_0\,\delta p_r-\lambda\,\Theta G_0]
+O(\varepsilon^2)$; integrating in $r$ gives equation~\eqref{eq:adiab-exact}. The reduction is
non-uniform at the turning point ($H_{p_r}\to0$): it holds on compact sub-arcs, and all
displayed residuals are taken there.

\paragraph{The universal boundary form of the Vaidya source.} The source does not
require a per-metric reduction to be written in closed form. Any arrival-time
Hamiltonian is Randers, $H=\beta(r)\,J+\sqrt{A_{rr}(r)\,p_r^2+A_{\varphi\varphi}(r)\,J^2}-1$,
homogeneous of degree one in $(p_r,J)$ apart from the $-1$, so Euler's theorem gives
\begin{equation}
    \big(J\,\partial_J+p_r\,\partial_{p_r}\big)H=H+1 .
    \label{eq:finsler-euler}
\end{equation}
Second, $H_v$ of~\eqref{eq:Hv} is invariant under the simultaneous rescaling
$(r,m,J)\mapsto(\alpha r,\alpha m,\alpha J)$ at fixed $p_r$---the Schwarzschild
self-similarity of the frozen family---so it is homogeneous of degree zero in $(r,m,J)$ and
Euler's theorem in those variables gives
\begin{equation}
    m\,\partial_m H=-\big(r\,\partial_r H+J\,\partial_J H\big)
    \qquad\text{identically, not merely on the shell.}
    \label{eq:selfsimilar}
\end{equation}
Combining the two turns the source integrand into a total derivative. Along the frozen
flow $\dd r/\dd\lambda=H_{p_r}$ and $\dd p_r/\dd\lambda=-H_r$, so
\[
   \frac{\dd (r\,p_r)}{\dd\lambda}=p_r H_{p_r}-r H_r
   \overset{\eqref{eq:finsler-euler}}{=}\big(H+1-J H_J\big)-r H_r
   \overset{\eqref{eq:selfsimilar}}{=}1+\Theta H+H ,
\]
and the shell condition $H=0$ removes the last term, leaving
$\Theta H=\dd(r\,p_r)/\dd\lambda-1$. Integrating from the terminal anchor gives the
\emph{universal boundary form}
\begin{equation}
    S_{\mathcal D}(\lambda)=[\,rp_r\,]_0^{\lambda}-\lambda,
    \label{eq:SD-vaidya}
\end{equation}
metric-independent up to the endpoint bracket. Identities~\eqref{eq:finsler-euler}
and~\eqref{eq:selfsimilar}, the expression~\eqref{eq:ThetaHv} for $\Theta H_v$, and the
combination $p_rH_{p_r}-rH_r-(1+\Theta H)=H$ were each verified symbolically in
\textsc{Mathematica}, all four residuals reducing to zero identically in
$(r,m,\Ehat,J,p_r)$; the numerical cross-check is
\texttt{VaidyaMetric/universal\_SD\_source\_vaidya.py}. The explicit Vaidya on-shell and
off-shell closed forms, with the horizon dilogarithm, are assembled in
section~\ref{sec:vaidya-adiabatic}.

\paragraph{The remainder is $O(\varepsilon^2)$.} The expansion above is an asymptotic
statement about the Hamiltonian flow, and on the sub-arcs where the reduction is used it
carries a genuine bound. Fix a compact sub-arc $[\lambda_-,\lambda_+]$ on which the frozen
solution stays in a compact subset $\mathcal K$ of the regular domain and on which
$|H_{p_r}|\ge\kappa>0$ (i.e.\ bounded away from the turning point). On a neighbourhood of
$\mathcal K$ the extended Hamiltonian $\mathcal H=p_s+H(r,p_r;m(s))$ is $C^2$ with bounded
derivatives, so writing $z=(r,p_r)$ and $z_\varepsilon$ for the exact flow, the exact and
first-order flows obey
\[
   \big|z_\varepsilon(\lambda)-z_0(\lambda)-\varepsilon\,z_1(\lambda)\big|
   \le C\,\varepsilon^{2}\,e^{L(\lambda-\lambda_-)} ,
   \qquad \lambda\in[\lambda_-,\lambda_+],
\]
by the standard Gr\"onwall estimate applied to the difference between the exact flow and
its first variation, with $L$ a bound for $\|\partial_z^2\mathcal H\|$ on $\mathcal K$ and
$C$ depending on the same $C^2$ bounds. Dividing by $H_{p_r}$, legitimate since
$|H_{p_r}|\ge\kappa$, transfers the estimate to $\dd\varphi/\dd r$ and, after integration
over the sub-arc, to $\delta\varphi$; this is the precise sense in
which~\eqref{eq:adiab-exact} closes the correction to $O(\varepsilon^2)$. The estimate
degrades as $\kappa\to0$, which is why the reduction is non-uniform at the turning point
and why all displayed residuals are taken on compact sub-arcs bounded away from it. The
measured log--log slope $\approx2.0$ of figure~\ref{fig:vaidya-offshell} is the numerical
confirmation of this exponent.

\paragraph{Special-function class.} The building blocks form a strict transcendence
hierarchy. The second-kind integrals $U_0,\dots,U_3=\int r^k/\sqrt S\,\dd r$ are
\emph{ordinary Abelian integrals} (weight one), reducible to the classical Kleinian
$\zeta,\sigma$ basis (Baker~\cite{BakerAbelian1897}; section~\ref{sec:vaidya-adiabatic}),
the same class as $\varphi_0$. The letters
$W_{jk}=\int(r^jU_k-r^kU_j)/\sqrt S\,\dd r$ are \emph{length-two iterated} Abelian
integrals with algebraic coefficients---one transcendence level above the $U_k$ and
\emph{not} reducible to the weight-one Kleinian class (their irreducibility to weight one
is \emph{conjectural}). The third-kind letter at the horizon $r=2m$---the advanced-time
$\dd\log$ letter, whose residue is nonzero, unlike the on-shell $\psi$ where it
vanishes---adds a genus-two dilogarithm (descriptively, a hyperelliptic polylogarithm
\cite{BauneEtal2024,DHokerHiddingSchlotterer2025}), evaluated by the geometrically
convergent theta-nome series in the reproducibility package; its irreducibility and
single-valued completion remain open---and the meromorphic sector in which this
letter sits is precisely where the higher-genus construction reports an obstruction,
simultaneous meromorphicity in two points being unattainable from genus two onwards
without additional marked points
\cite[\S4.5]{DHokerHiddingSchlotterer2025}. Thus the \emph{frozen} reduction
$\partial_\lambda\varphi_0$ closes in the weight-one Abelian/Kleinian class, but already
the \emph{on-shell} first-order response carries a genuine weight-two sector---the
length-two iterated integrals $W_{jk}$; the off-shell term adds further length-two letters
and one genus-two (horizon) dilogarithm. Reducibility of the $W_{jk}$ to weight one and a
canonical single-valued completion remain conjectural.

\paragraph{What is reconstructible from the text alone.} The manuscript fixes the
algebraic reduction and the closed special-function representation without the
repository: the frozen integrands
$\dd\varphi/\dd r=\mathcal K/\sqrt S$ with their spectral radicals and numerators;
the on-shell reduction (the polynomial identity $2A'S-AS'+2S\sum_k c_kr^k=2N_\lambda$
and its $11\times11$ linear system---no fit); all block coefficients $c_k,g_i,Q_{kj}$;
the Kleinian $\zeta,\sigma$ closed forms of the $U_k$ (Hermite recursion) and the assembly
of the weight-two $W_{jk}$ from them (not reduced to the weight-one class); the elementary
remainders; and the defining \emph{form} of the genus-two letter
$L=\log[\theta[\delta](w-e_+)/\theta[\delta](w-e_-)]$ with its odd half-characteristic
$\delta$. The text does \emph{not} print a standalone nome-series reconstruction recipe.
Numerical evaluation of the genus-two $\theta$---constructing the period matrix
$\tau$ from $Y^2=\mathcal S$, its Siegel reduction, and summing the geometrically
convergent nome series (order $N{=}4$ already to $10^{-13}$) are performed by
\texttt{Sage}$+$\texttt{abelfunctions} (\texttt{RiemannTheta}) in the reproducibility
package~\cite{Rosignoli2026CodeJMP}. In short: the PDF is self-contained at the level of the
algebraic reduction and closed special-function representation; period construction and
genus-two numerical values require the archived routine.

\paragraph{Branch forms.} For the \emph{Vaidya} $v$-branch the slow datum $m$ is a metric
modulus, \emph{not} a conformal rescaling of momenta: there is no $p_r/A$ transformation
and hence no dilation term, so $\mathcal D=\Theta=m\partial_m$ is complete and
\begin{equation}
    \Theta H_v=m\,\partial_m H_v
    =m\!\left[-\frac{2p_r}{r}+\frac{\sqrt{\Ehat^2p_r^2+J^2/r^2}}{r\sqrt w}\right],
    \qquad w=\Ehat^2-f;
    \label{eq:ThetaHv}
\end{equation}
since $m$ \emph{increases} along the ingoing flow while the elapsed clock $\lambda$ is
oriented inward from the free-arrival endpoint, $\lambda$ runs \emph{opposite} to the mass
drift on this branch. This fixed orientation ($\mathrm{sgn}(\dd\lambda/\dd m)=-1$) is what
inverts both source contributions in equation~\eqref{eq:adiab-exact}, i.e.\
$\lambda\Theta H_0-S_{\mathcal D}\to S_{\mathcal D}-\lambda\Theta H_0$ and
$-\lambda\Theta G_0\to+\lambda\Theta G_0$---a declared sign convention, not a change of
result. The proper-time branch uses $H_\tau$ of equation~\eqref{eq:Htau-vaidya}. Both Vaidya branches (advanced-time $v$ and proper-time
$\tau$) close to $O(\varepsilon^2)$ against the true flow: least-squares log--log slope
$\approx1$ for the on-shell term alone versus $\approx2.0$ for the complete first-order
term (figure~\ref{fig:vaidya-offshell}).

\paragraph{Algebraic genus-degeneration limit ($J\to\Jc$).} The degeneration form is
\emph{not} obtained by substituting $J=\Jc$ into the generic expressions. The source
\emph{numerator} $N_\lambda$ is polynomial in $J$ and hence perfectly finite at $\Jc$, but
the reduction \emph{coefficients} $c_k$ are not: by~\eqref{eq:ck-closed} they carry the
reduced discriminant $\Delta$ in the denominator, and $\Delta$ vanishes precisely at
$|J|=\Jc$, so each $c_k$ has a simple pole there---the analytic counterpart of the rank
drop of Lemma~\ref{lem:reduction}. The degeneration forms below are therefore constructed
\emph{directly} on the degenerate curve, $S=(r-r_d)^2Q_4$ (double root $r_d$, genus two
$\to$ genus one), and not as a limit of the genus-two formulas. Each hyperelliptic letter collapses,
$U_k=\int r^k/\sqrt S\to V_k+r_d^{\,k}\,\Pi_{r_d}$ with $V_k=\int(\text{poly})/\sqrt{Q_4}$
(elliptic second kind, Weierstrass $\wp,\zeta$) and $\Pi_{r_d}=\int\dd r/[(r-r_d)\sqrt{Q_4}]$
(third kind); the spurious poles at $r_d$ cancel in the physical combination (as
$S(r_d)=0$), leaving the weight-two block as a genuine \emph{elliptic}
dilogarithm~\cite{BeilinsonLevin1994,BrownLevin2011} plus the horizon third kind. Thus the
degeneration off-shell form is one degree \emph{lower} than the generic genus-two case;
verified letter by letter to machine precision
(\texttt{VaidyaMetric/vaidya\_sep\_offshell\_elliptic.py}). For Vaidya the double root is a
branch point \emph{off} the physical arc and $\mathcal D=\Theta$ carries no dilation, so
the degeneration limit is regular and elliptic (non-uniform only at $r_d$) and no
orbit crosses it under the mass drift. The rotating case is qualitatively
different---there the conformal dilation ray drives the orbit \emph{across} the moving
separatrix, a Neishtadt separatrix-crossing problem
\cite{Neishtadt1986} treated in Appendix~C of Paper~II.



\FloatBarrier


% ---- back matter -----------------------------------------------------
\section*{Author declarations}

\subsection*{Conflict of interest}
The author has no conflicts to disclose.

\subsection*{Author contributions}
\textbf{Iman Rosignoli}: Conceptualization (lead); Formal analysis (lead); Investigation
(lead); Methodology (lead); Software (lead); Validation (lead); Visualization (lead);
Writing --- original draft (lead); Writing --- review \& editing (lead).

\subsection*{Ethics approval}
Not applicable: this study involves no human participants and no animals.

\subsection*{Funding}
This research received no external funding.

\section*{Data availability}
The data that support the findings of this study are openly available in the Zenodo
repository \emph{NonStationaryBrachistochrone}, version~1.2.0, at
\url{https://doi.org/10.5281/zenodo.22035415}, reference~\cite{Rosignoli2026CodeJMP}. This
comprises the numerical data together with the symbolic and numerical code, released under
the MIT licence; the development history is available in the associated GitHub repository
\url{https://github.com/ImanetorX98/NonStationaryBrachistochrone}.

\ack
The author thanks the University of Pavia for a stimulating environment. OpenAI
GPT-5.6 (Sol) and Anthropic Claude (Opus~4.8) were used to assist with symbolic-algebra
checks, numerical-code development, and language and structural editing of the manuscript;
all AI-assisted output was critically reviewed and independently verified by the author,
who takes full responsibility for the accuracy, originality and integrity of the work. The
closed-form and hyperelliptic computations relied on \textsc{SageMath}~\cite{SageMath}
with \textsc{abelfunctions}~\cite{abelfunctions}, and on
\textsc{SymPy}~\cite{SymPy}, \textsc{NumPy}~\cite{NumPy}, \textsc{SciPy}~\cite{SciPy} and
\textsc{Matplotlib}~\cite{Matplotlib}.

\medskip\noindent\textit{ORCID iD:} Iman Rosignoli,
\url{https://orcid.org/0009-0004-4536-0285}.

\section*{References}
\bibliographystyle{iopart-num}
\bibliography{../paper/refs,../paper/companionII}

\end{document}
